% -----------------------------------------------------------
\section{1838}
% -----------------------------------------------------------

	% ----------------------
	\subsection{Joseph Smith}
	% ----------------------
		\label{EMS-1838-JS}

		% ----------------------
		\subsubsection{Introduction to 1838}
		% ----------------------
		
			sdf
			
		% -----
		\subsubsection{Letter and Revelation to Edward Partridge, 7 January 1838}
		% -----
			\label{EMS-1838-JS-7-JAN-1838}
			% \label{EMS-1830-JS-DAY-3LETTERMONTH-YEAR}
			
			\begin{description}
					\item[Print Sources]
				\end{description}

					\begin{tabular}{p{0.5in}p{1in}p{0.5in}p{1in}}
					JSP	 	& D5:492--494		& JSR 		& 284 \\
					DC	 	& ---		& HC 		& --- \\
					HDDC 	& ---		& TCHC 		& --- \\
					\end{tabular}
					% HDDC = woodford's DC diss; JSR = Marquardt's JS Revelations; TCHC = Vogel HC
				
				\begin{description}
					\item[Online Access] \href{https://www.josephsmithpapers.org/paper-summary/letter-and-revelation-to-edward-partridge-7-january-1838/1}{Here}
					% \item[Analysis] \hyperref[XX]{Here}
				\end{description}
				
				\begin{description}
					\item[Background]
				\end{description}
				
					% It might also be useful to give a two sentence context of the period: e.g., "This speech was given in Nauvoo to a small congregation one month prior to the destruction of the Nauvoo Expositor. These and other tensions may have been on the prophet's mind when he gave the speech."
				
				\begin{description}
					\item[Original Source]
				\end{description}
				
					% brief description of original source material, with reference to JSP or somewhere else for more info
					% Background material: it might be helpful to have a note that says whether a given document was presented as a speech, was written in a journal, etc.
					% Also a comment here about how reliable the source might be, or what shape it is in, if there are large omissions or parts that are hard to understand, and so on.
				
				\begin{description}
					\item[Text]
				\end{description}
					
					\begin{tightcenter}
					Kirtland, Jan. 7\supersu{th} 1838
					\end{tightcenter}
					
					Brother [Edward] Partridge, Thus saith the Lord, my servant Edward and his house shall be numbered with the blessed and Abraham their father, and his name shall be had in sacred rem[em]brance. And again thus saith the Lord, let my people be aware of dissensions among them lest the enemy have power <over> them, Awake my shepherds and warn my people! for behold the wolf cometh to destroy them! receive him not. Now I would inform you that my health is good and also that of my family, and it is my earnest prayer to God that health, peace, and pleanty may crown your board and blessings of Heaven rest upon the head of him in whom the Lord hath said there is no guile.
					
		% -----
		\subsubsection{Revelation, 12 January 1838--A}
		% -----
			\label{EMS-1838-JS-12-JAN-1838-A}
			% \label{EMS-1830-JS-DAY-3LETTERMONTH-YEAR}
			
			\begin{description}
					\item[Print Sources]
				\end{description}

					\begin{tabular}{p{0.5in}p{1in}p{0.5in}p{1in}}
					JSP	 	& D5:494--498; J1:281--283		& JSR 		& 284--285 \\
					DC	 	& ---		& HC 		& --- \\
					HDDC 	& ---		& TCHC 		& --- \\
					\end{tabular}
					% HDDC = woodford's DC diss; JSR = Marquardt's JS Revelations; TCHC = Vogel HC
				
				\begin{description}
					\item[Online Access] \href{https://www.josephsmithpapers.org/paper-summary/revelation-12-january-1838-a/1}{Here}
					% \item[Analysis] \hyperref[XX]{Here}
				\end{description}
				
				\begin{description}
					\item[Background]
				\end{description}
				
					% It might also be useful to give a two sentence context of the period: e.g., "This speech was given in Nauvoo to a small congregation one month prior to the destruction of the Nauvoo Expositor. These and other tensions may have been on the prophet's mind when he gave the speech."
				
				\begin{description}
					\item[Original Source]
				\end{description}
				
					% brief description of original source material, with reference to JSP or somewhere else for more info
					% Background material: it might be helpful to have a note that says whether a given document was presented as a speech, was written in a journal, etc.
					% Also a comment here about how reliable the source might be, or what shape it is in, if there are large omissions or parts that are hard to understand, and so on.
				
				\begin{description}
					\item[Text]
				\end{description}
				
					\begin{tightcenter}
						\hypertarget{EMS-1838-JS-12-JAN-1838-A-a}%
					Kirtland
						\marginnote{a}%
					Ohio Ja\uline{n.} 12\supers{th} 1838
					\end{tightcenter}

					In the presence of Joseph Smith Jr Sidney Rigdon Vinson Knight \& G[eorge] W. Robinson at the French Farm, The following inquiry was made of the Lord

					A question asked of the Lord concerning the trying of the first Presidency of the Church of Latter Day Saints for transgression according to the item of Law found in third sec of the Book of covenants 37 verse Whether the decision of such an council of one stake shall be conclusive for Zion and all the Stake------

					Answer\} Thus saith the Lord the time has now come when a decision of such an council would not answer for Zion and all her stakes------

						\hypertarget{EMS-1838-JS-12-JAN-1838-A-b}%
					What
						\marginnote{b}%
					will answer for Zion and all her Stakes

					Answer; Thus saith the Lord let the first presidency of my Church be held in full fellowship in Zion and all her stakes until they shall be found transgressors by such an high council------ as is named in the 3\supers{rd} sec. 37 verse of the Book of Covenants, in Zion by 3 witnesses standing against Each member of said presidency and said witnesses shall be of long and faithful standing and such also as cannot be impeached by other witnesses before said council and when a desision is had by such an council in Zion it shall only be for Zion it shall not answer for her stakes but if said desision be acknowleged by the Council of her stakes then it shall answer for her stakes <But if it is not acknowleged by the Stakes then such <stakes> may have the privilege of hearing for themselves> or if said decision shall be acknowleged by a majority of her stakes then it shall <answer> for <all> her stakes
						\hypertarget{EMS-1838-JS-12-JAN-1838-A-c}%
					And
						\marginnote{c}%
					again the presidency of said Church may be tried by the <voice> of the whole body of <the Church of> Zion, and the \sout{whole body} <voice of a majority> of all her stakes, And again except a majority is had by the voice of the Church of Zion, and the majority of her stakes, the charges will be considered not sustained, and in order to sustain such charge or charges before said Church <of Zion> or her stakes Such witnesses must be had as is named above, That is three witnesses $\diamond$$\diamond$ Each president that is of long \& faithfull sta[nd]ing that can[not] be impeached by other wi[tnesses] before [the] Church of Zion or her stakes, and all this saith the Lord, because of wicked and aspiring men, let all your doing be in meekness and humility before me even so Amen

					[\textit{page blank}]

		% -----
		\subsubsection{Revelation, 12 January 1838--B}
		% -----
			\label{EMS-1838-JS-12-JAN-1838-B}
			% \label{EMS-1830-JS-DAY-3LETTERMONTH-YEAR}
			
			\begin{description}
					\item[Print Sources]
				\end{description}

					\begin{tabular}{p{0.5in}p{1in}p{0.5in}p{1in}}
					JSP	 	& D5:498--499; J1:283		& JSR 		& 285 \\
					DC	 	& ---		& HC 		& --- \\
					HDDC 	& ---		& TCHC 		& --- \\
					\end{tabular}
					% HDDC = woodford's DC diss; JSR = Marquardt's JS Revelations; TCHC = Vogel HC
				
				\begin{description}
					\item[Online Access] \href{https://www.josephsmithpapers.org/paper-summary/revelation-12-january-1838-b/1}{Here}
					% \item[Analysis] \hyperref[XX]{Here}
				\end{description}
				
				\begin{description}
					\item[Background]
				\end{description}
				
					% It might also be useful to give a two sentence context of the period: e.g., "This speech was given in Nauvoo to a small congregation one month prior to the destruction of the Nauvoo Expositor. These and other tensions may have been on the prophet's mind when he gave the speech."
				
				\begin{description}
					\item[Original Source]
				\end{description}
				
					% brief description of original source material, with reference to JSP or somewhere else for more info
					% Background material: it might be helpful to have a note that says whether a given document was presented as a speech, was written in a journal, etc.
					% Also a comment here about how reliable the source might be, or what shape it is in, if there are large omissions or parts that are hard to understand, and so on.
				
				\begin{description}
					\item[Text]
				\end{description}
				
					\begin{tightcenter}
					Kirtland Jan 12\supers{th} 1838
					\end{tightcenter}
					
					Can \sout{the first} any branch of the Church of Latter Day Saints be concidered a stake of Zion utill [until] they have acknowleged the authority of the first Presidensy, by a vote of said Church,
					
					Thus saith the Lord verly I say unto you nay
					
					How then
					
					Answer. No stake shall be appointed except by the first presidency and this Presidency be acknowleged by the voice of the same, otherwise it shall not be counted as a stake of Zion, And again except it be dedicated by this Presidency it cannot be acknowleged as a stake of Zion, For unto this end, have I appointed them, in laying the foundation of and establishing my Kingdom
					
		% -----
		\subsubsection{Revelation, 12 January 1838--C}
		% -----
			\label{EMS-1838-JS-12-JAN-1838-C}
			% \label{EMS-1830-JS-DAY-3LETTERMONTH-YEAR}
			
			\begin{description}
					\item[Print Sources]
				\end{description}

					\begin{tabular}{p{0.5in}p{1in}p{0.5in}p{1in}}
					JSP	 	& D5:500--502; J1:283--284		& JSR 		& 286 \\
					DC	 	& ---		& HC 		& --- \\
					HDDC 	& ---		& TCHC 		& --- \\
					\end{tabular}
					% HDDC = woodford's DC diss; JSR = Marquardt's JS Revelations; TCHC = Vogel HC
				
				\begin{description}
					\item[Online Access] \href{https://www.josephsmithpapers.org/paper-summary/revelation-12-january-1838-c/1}{Here}
					% \item[Analysis] \hyperref[XX]{Here}
				\end{description}
				
				\begin{description}
					\item[Background]
				\end{description}
				
					% It might also be useful to give a two sentence context of the period: e.g., "This speech was given in Nauvoo to a small congregation one month prior to the destruction of the Nauvoo Expositor. These and other tensions may have been on the prophet's mind when he gave the speech."
				
				\begin{description}
					\item[Original Source]
				\end{description}
				
					% brief description of original source material, with reference to JSP or somewhere else for more info
					% Background material: it might be helpful to have a note that says whether a given document was presented as a speech, was written in a journal, etc.
					% Also a comment here about how reliable the source might be, or what shape it is in, if there are large omissions or parts that are hard to understand, and so on.
				
				\begin{description}
					\item[Text]
				\end{description}
				
					Thus saith the Lord Let the presidency of my Church take their families as soon as it is praticable and a door is open for them and move on to the west as fast as the way is made plain before their faces and let their hearts be comforted for I will be with them
					
					Verily I say unto you the time [has] come that your laibours are finished in this place, for a season, Therefore arise and get yourselves on to a land which I shall show unto you even a land flowing with milk and honey you are clean from the blood of this people and wo unto those who have become your enimies who <have> professed my name saith the Lord, for their judgement lingereth not and their damnation slumbereth not, let all your faithfull friends arise with their families also and get out of this place and gather themselves together unto Zion and be at peace among yourselves O ye inhabitants of Zion or there shall be no saf[e]ty for you
					
					[\textit{page blank}]
					
					[\textit{page blank}]
					
		% -----
		\subsubsection{Minutes, 15 March 1838}
		% -----
			\label{EMS-1838-JS-15-MAR-1838}
			% \label{EMS-1830-JS-DAY-3LETTERMONTH-YEAR}
			
			\begin{description}
					\item[Print Sources]
				\end{description}

					\begin{tabular}{p{0.5in}p{1in}p{0.5in}p{1in}}
					JSP	 	& D6:39--43		& JSR 		& --- \\
					DC	 	& ---		& HC 		& --- \\
					HDDC 	& ---		& TCHC 		& --- \\
					\end{tabular}
					% HDDC = woodford's DC diss; JSR = Marquardt's JS Revelations; TCHC = Vogel HC
				
				\begin{description}
					\item[Online Access] \href{https://www.josephsmithpapers.org/paper-summary/minutes-15-march-1838/1}{Here}
					% \item[Analysis] \hyperref[XX]{Here}
				\end{description}
				
				\begin{description}
					\item[Background]
				\end{description}
				
					% It might also be useful to give a two sentence context of the period: e.g., "This speech was given in Nauvoo to a small congregation one month prior to the destruction of the Nauvoo Expositor. These and other tensions may have been on the prophet's mind when he gave the speech."
				
				\begin{description}
					\item[Original Source]
				\end{description}
				
					% brief description of original source material, with reference to JSP or somewhere else for more info
					% Background material: it might be helpful to have a note that says whether a given document was presented as a speech, was written in a journal, etc.
					% Also a comment here about how reliable the source might be, or what shape it is in, if there are large omissions or parts that are hard to understand, and so on.
				
				\begin{description}
					\item[Text]
				\end{description}
				
					President Joseph Smith j\supersu{r} gave a history of the ordination of David Whitmer, which took place in July 1834, to be a leader, or a prophet to this Church, which (ordination) was on conditions that he (J. Smith jr) did not live to God himself.
					
					President J. Smith jr approved of the proceedings of the High Council, after hearing the minutes of the former Councils.
					
		% -----
		\subsubsection{Motto, circa 16 or 17 March 1838}
		% -----
			\label{EMS-1838-JS-CIR-16-17-MAR-1838}
			% \label{EMS-1830-JS-DAY-3LETTERMONTH-YEAR}
			
			\begin{description}
					\item[Print Sources]
				\end{description}

					\begin{tabular}{p{0.5in}p{1in}p{0.5in}p{1in}}
					JSP	 	& D6:43--45; J1:237--238		& JSR 		& --- \\
					DC	 	& ---		& HC 		& 3:9 \\
					HDDC 	& ---		& TCHC 		& 3:11 \\
					TPJS 	& 117			& 	 		&  \\
					\end{tabular}
					% HDDC = woodford's DC diss; JSR = Marquardt's JS Revelations; TCHC = Vogel HC
				
				\begin{description}
					\item[Online Access] \href{https://www.josephsmithpapers.org/paper-summary/motto-circa-16-or-17-march-1838/1}{Here}
					% \item[Analysis] \hyperref[XX]{Here}
				\end{description}
				
				\begin{description}
					\item[Background]
				\end{description}
				
					% It might also be useful to give a two sentence context of the period: e.g., "This speech was given in Nauvoo to a small congregation one month prior to the destruction of the Nauvoo Expositor. These and other tensions may have been on the prophet's mind when he gave the speech."
				
				\begin{description}
					\item[Original Source]
				\end{description}
				
					% brief description of original source material, with reference to JSP or somewhere else for more info
					% Background material: it might be helpful to have a note that says whether a given document was presented as a speech, was written in a journal, etc.
					% Also a comment here about how reliable the source might be, or what shape it is in, if there are large omissions or parts that are hard to understand, and so on.
				
				\begin{description}
					\item[Text]
				\end{description}
				
					Motto of the Church of Christ of \uline{Latterday} \uline{Saints}.

					The Constitution of our country formed by the Fathers of Liberty.

					Peace and good order in society Love to God and good will to man.

					All good and wholesome Law's; And virtue and truth above all things

					And Aristarchy live forever!!!

					But Wo to tyrants, Mobs, Aristocracy, Anarchy and Toryism: And all those who invent or seek out unrighteous and vexatious lawsuits under the pretext or color of law or office, either religious or political.
					
					Exalt the standard of Democracy! Down with that of Priestcraft, and let all the people say Amen! that the blood of our Fathers may not cry from the ground against us.

					Sacred is the Memory of that Blood which baught for us our liberty.

					\begin{tabular}{p{3cm}p{3cm}p{4cm}p{4cm}}
					Signed &  &  & Joseph Smith J\supersu{r}\supers{.} \\
					 & \uline{Geo. W. Robinson} &  & Thomas B. Marsh \\
					 &  &  & D[avid] W. Patten \\
					 &  &  & Brigham Youngs [Young] \\
					 &  &  & Samuel H. Smith \\
					 &  &  & George M. Hinkle \\
					 &  &  & \uline{John Corrill}.---
					\end{tabular}
					
		% -----
		\subsubsection{Questions and Answers, between circa 16 and circa 29 March 1838--A [D\&C 113:1--6]}
		% -----
			\label{EMS-1838-JS-CIR-16-CIR-29-MAR-1838-A}
			% \label{EMS-1830-JS-DAY-3LETTERMONTH-YEAR}
			
			\begin{description}
					\item[Print Sources]
				\end{description}

					\begin{tabular}{p{0.5in}p{1in}p{0.5in}p{1in}}
					JSP	 	& D6:50--54				& JSR 		& 286--287 \\
					DC	 	& Section 113		& HC 		& 3:9--10 \\
					HDDC 	& 1495--1499		& TCHC 		& 3:12 \\
					\end{tabular}
					% HDDC = woodford's DC diss; JSR = Marquardt's JS Revelations; TCHC = Vogel HC
				
				\begin{description}
					\item[Online Access] \href{https://www.josephsmithpapers.org/paper-summary/questions-and-answers-between-circa-16-and-circa-29-march-1838-a-dc-1131-6/1}{Here}
					% \item[Analysis] \hyperref[XX]{Here}
				\end{description}
				
				\begin{description}
					\item[Background]
				\end{description}
				
					% It might also be useful to give a two sentence context of the period: e.g., "This speech was given in Nauvoo to a small congregation one month prior to the destruction of the Nauvoo Expositor. These and other tensions may have been on the prophet's mind when he gave the speech."
				
				\begin{description}
					\item[Original Source]
				\end{description}
				
					% brief description of original source material, with reference to JSP or somewhere else for more info
					% Background material: it might be helpful to have a note that says whether a given document was presented as a speech, was written in a journal, etc.
					% Also a comment here about how reliable the source might be, or what shape it is in, if there are large omissions or parts that are hard to understand, and so on.
				
				\begin{description}
					\item[Text]
				\end{description}
				
					Quest. on Scripture.

					1\supersu{st}\supers{.} Who is the stem of Jessee spoken of in the 1\supersu{st}\supers{.} 2\supersu{d}\supers{.} 3\supersu{d}\supers{.} 4\supersu{th}\supers{.} and 5\supersu{th}\supers{.} verses of the 11\supersu{th}\supers{.} Chap. of Isiah.

					Ans. Verely thus saith the Lord It is Christ

					Q. 2\supersu{d}\supers{.} What is the Rod spoken of in the 1\supersu{st}\supers{.} verse of the 11\supersu{th}\supers{.} \sout{verse} Chap. that shoud come of the stem of Jessee.

					Ans. Behold thus saith <the Lord> it is a servant in the hands of Christ who is partly a decendant of Jessee as well as of Ephraim or of the house of Joseph, on whome thare is Laid much power.

					Qest 3\supersu{d}\supers{.} What is the Root of Jessee spoken of in the 10\supersu{th}\supers{.} verse of the 11\supersu{th}\supers{.} Chap.

					Ans. Behold thus saith the Lord; it is a decendant of Jessee as well as of Joseph unto whom rightly belongs the Priesthood and the kees of the Kingdom for an ensign and for the geathering of my people in the Last day.---
					
		% -----
		\subsubsection{Questions and Answers, between circa 16 and circa 29 March 1838--B [D\&C 113:7--10]}
		% -----
			\label{EMS-1838-JS-CIR-16-CIR-29-MAR-1838-B}
			% \label{EMS-1830-JS-DAY-3LETTERMONTH-YEAR}
			
			\begin{description}
					\item[Print Sources]
				\end{description}

					\begin{tabular}{p{0.5in}p{1in}p{0.5in}p{1in}}
					JSP	 	& D6:54--56				& JSR 		& 286--287 \\
					DC	 	& Section 113		& HC 		& 3:9--10 \\
					HDDC 	& 1495--1499		& TCHC 		& 3:12 \\
					\end{tabular}
					% HDDC = woodford's DC diss; JSR = Marquardt's JS Revelations; TCHC = Vogel HC
				
				\begin{description}
					\item[Online Access] \href{https://www.josephsmithpapers.org/paper-summary/questions-and-answers-between-circa-16-and-circa-29-march-1838-b-dc-1137-10/1}{Here}
					% \item[Analysis] \hyperref[XX]{Here}
				\end{description}
				
				\begin{description}
					\item[Background]
				\end{description}
				
					% It might also be useful to give a two sentence context of the period: e.g., "This speech was given in Nauvoo to a small congregation one month prior to the destruction of the Nauvoo Expositor. These and other tensions may have been on the prophet's mind when he gave the speech."
				
				\begin{description}
					\item[Original Source]
				\end{description}
				
					% brief description of original source material, with reference to JSP or somewhere else for more info
					% Background material: it might be helpful to have a note that says whether a given document was presented as a speech, was written in a journal, etc.
					% Also a comment here about how reliable the source might be, or what shape it is in, if there are large omissions or parts that are hard to understand, and so on.
				
				\begin{description}
					\item[Text]
				\end{description}
				
					\uline{Questions} by \uline{Elias} \uline{Higby} [Higbee]

					1\supersu{st}\supers{.} Q. What is ment by the command in Isiah 52\supersu{d}\supers{.} chap 1\supersu{st}\supers{.} verse which saith Put on thy strength O Zion and what people had I[sa]iah referance to

					A. He had reference to those whome God should call in the last day's who should hold the power of Priesthood to bring again zion and the redemption of Israel.

					And to put on her strength is to put on the authority of the priesthood which she (zion) has a right to by lineage: Also to return to that power which she had lost

					Ques. 2\supersu{d}\supers{.} What are we to understand by zions loosing herself from the bands of her neck 2\supersu{d}\supers{.} verse.

					A. We are to understand that the scattered remnants are exorted to to return to the Lord from whence they have fal[l]en which if they do the promise of the Lord is that he will speak to them or give them revelation See 6\supersu{th}\supers{.} 7\supersu{th}\supers{.} and 8\supersu{th}\supers{.} verses The bands of her neck are the curses of God upon her or the remnants of Israel in their scattered condition among the Gentiles.
					
		% -----
		\subsubsection{Letter to the Presidency in Kirtland, 29 March 1838}
		% -----
			\label{EMS-1838-JS-29-MAR-1838}
			% \label{EMS-1830-JS-DAY-3LETTERMONTH-YEAR}
			
			\begin{description}
					\item[Print Sources]
				\end{description}

					\begin{tabular}{p{0.5in}p{1in}p{0.5in}p{1in}}
					JSP	 	& D6:56--61; J1:245--247		& JSR 		& --- \\
					DC	 	& ---		& HC 		& 3:10--12 \\
					HDDC 	& ---		& TCHC 		& 3:12--14 \\
					\end{tabular}
					% HDDC = woodford's DC diss; JSR = Marquardt's JS Revelations; TCHC = Vogel HC
				
				\begin{description}
					\item[Online Access] \href{https://www.josephsmithpapers.org/paper-summary/letter-to-the-presidency-in-kirtland-29-march-1838/1}{Here}
					% \item[Analysis] \hyperref[XX]{Here}
				\end{description}
				
				\begin{description}
					\item[Background]
				\end{description}
				
					% It might also be useful to give a two sentence context of the period: e.g., "This speech was given in Nauvoo to a small congregation one month prior to the destruction of the Nauvoo Expositor. These and other tensions may have been on the prophet's mind when he gave the speech."
				
				\begin{description}
					\item[Original Source]
				\end{description}
				
					% brief description of original source material, with reference to JSP or somewhere else for more info
					% Background material: it might be helpful to have a note that says whether a given document was presented as a speech, was written in a journal, etc.
					% Also a comment here about how reliable the source might be, or what shape it is in, if there are large omissions or parts that are hard to understand, and so on.
				
				\begin{description}
					\item[Text]
				\end{description}
				
						\hypertarget{EMS-1838-JS-29-MAR-1838-a}%
					...Being
						\marginnote{a}%
					under the hand of wicked vexatious Lawsuits for seven years past my buisness was so dangerous that I was not able to leave it, in as good a situation as I had antisipated, but if there are any wrongs, They shall all be noticed so far as the Lord gives me ability \& power to do so, say to all the brotheren that I have not forgotten them, but remember them in my prayers, Say to Mother Beaman [Sarah Burt Beman] that I remembr her, Also Br Daniel Carter Br Stong \& family Br [Oliver] Granger \& family, Finally I cannot innumerate them all for the want of room I will just name Br Knights the Bishop \&c. My best respects to them all \sout{for the want of room} \& I commend them and the Church of God in Kirtland to our Heavenly Father \& the word of his grace, which is able to make you wise unto Salvation
						\hypertarget{EMS-1838-JS-29-MAR-1838-b}%
					I
						\marginnote{b}%
					would just say to Br. [William] Marks, that I saw in a vision while on the road that whereas he was closely persued by an innumerable concource of enimies and as they pressed upon him hard as if they were about to devour him, \sout{It} <\&> had seemingly attained some degre[e] of advantage over him But about this time a chariot of fire came and near the place and the Angel of the Lord put forth his hand unto Br. Marks \& said unto him thou art my son come \uline{here}, and immediately he was caught up in the Chariot and rode away triumphantly out of their midst and again the Lord said I will raise th[ee] up for a blessing unto many people Now the particulars of this whole matter cannot be writen at this time but the vision was evidently given to me that I might know that the hand of the Lord would be on his behalf...
					
		% -----
		\subsubsection{Minutes, 7--8 April 1838}
		% -----
			\label{EMS-1838-JS-7-8-APR-1838}
			% \label{EMS-1830-JS-DAY-3LETTERMONTH-YEAR}
			
			\begin{description}
					\item[Print Sources]
				\end{description}

					\begin{tabular}{p{0.5in}p{1in}p{0.5in}p{1in}}
					JSP	 	& D6:70--74		& JSR 		& --- \\
					DC	 	& ---		& HC 		& 3:14--15 \\
					HDDC 	& ---		& TCHC 		& 3:16--18 \\
					TPJS 	& 117			& 	 		&  \\
					\end{tabular}
					% HDDC = woodford's DC diss; JSR = Marquardt's JS Revelations; TCHC = Vogel HC
				
				\begin{description}
					\item[Online Access] \href{https://www.josephsmithpapers.org/paper-summary/minutes-7-8-april-1838/1}{Here}
					% \item[Analysis] \hyperref[XX]{Here}
				\end{description}
				
				\begin{description}
					\item[Background]
				\end{description}
				
					% It might also be useful to give a two sentence context of the period: e.g., "This speech was given in Nauvoo to a small congregation one month prior to the destruction of the Nauvoo Expositor. These and other tensions may have been on the prophet's mind when he gave the speech."
				
				\begin{description}
					\item[Original Source]
				\end{description}
				
					% brief description of original source material, with reference to JSP or somewhere else for more info
					% Background material: it might be helpful to have a note that says whether a given document was presented as a speech, was written in a journal, etc.
					% Also a comment here about how reliable the source might be, or what shape it is in, if there are large omissions or parts that are hard to understand, and so on.
				
				\begin{description}
					\item[Text]
				\end{description}
				
					...President J. Smith jr. made some remarks, also gave some instruction respecting the order of the day. The conference was then opened by singing, ``O God our hope in ages past'' and prayer by President B. Young.--- Also a hymn was sung ``how firm a foundation''. After which, President J. Smith, Jr. arose and addressed the congregation at considerable length, on some important items...
					
					...President Joseph Smith Jr. made a few remarks respecting the Kirtland Bank--- Who was followed by Brigham Young, who gave a short history of his travels to Massachusetts and New York...
					
					...President Joseph Smith Jr, next made a few remarks on the word of wisdom, giving the reason of its coming forth, saying it should be observed. On motion, the Conference adjourned for one hour...
					
		% -----
		\subsubsection{Revelation, 11 April 1838 [D\&C 114]}
		% -----
			\label{EMS-1838-JS-11-APR-1838}
			% \label{EMS-1830-JS-DAY-3LETTERMONTH-YEAR}
			
			\begin{description}
					\item[Print Sources]
				\end{description}

					\begin{tabular}{p{0.5in}p{1in}p{0.5in}p{1in}}
					JSP	 	& D6:81--82; J1:257				& JSR 		& 287 \\
					DC	 	& Section 114		& HC 		& 3:23 \\
					HDDC 	& 1500--1505		& TCHC 		& 3:23 \\
					\end{tabular}
					% HDDC = woodford's DC diss; JSR = Marquardt's JS Revelations; TCHC = Vogel HC
				
				\begin{description}
					\item[Online Access] \href{https://www.josephsmithpapers.org/paper-summary/revelation-11-april-1838-dc-114/1}{Here}
					% \item[Analysis] \hyperref[DC114]{Here}
				\end{description}
				
				\begin{description}
					\item[Background]
				\end{description}
				
					% It might also be useful to give a two sentence context of the period: e.g., "This speech was given in Nauvoo to a small congregation one month prior to the destruction of the Nauvoo Expositor. These and other tensions may have been on the prophet's mind when he gave the speech."
				
				\begin{description}
					\item[Original Source]
				\end{description}
				
					% brief description of original source material, with reference to JSP or somewhere else for more info
					% Background material: it might be helpful to have a note that says whether a given document was presented as a speech, was written in a journal, etc.
					% Also a comment here about how reliable the source might be, or what shape it is in, if there are large omissions or parts that are hard to understand, and so on.
				
				\begin{description}
					\item[Text]
				\end{description}
				
					Revelation to D[avid] W. Patten. given April 11\supersu{th}\supers{.} 1838 Verily thus Saith the Lord, it is wisdom in my Servant D. W. Patten, that he settle up all his buisness, as soon as he possibly, can, and make a disposition of his merchandise, that he may perform a mission unto me next spring, in company with others even twelve including himself, to testify of my name and bear glad tidings unto all the world, for verrily thus Saith the Lord that inasmuch as there are those among you who deny my name, others shall be planted in their stead and receive their bishoprick Amen.------
					
		% -----
		\subsubsection{Revelation, 17 April 1838}
		% -----
			\label{EMS-1838-JS-17-APR-1838}
			% \label{EMS-1830-JS-DAY-3LETTERMONTH-YEAR}
			
			\begin{description}
					\item[Print Sources]
				\end{description}

					\begin{tabular}{p{0.5in}p{1in}p{0.5in}p{1in}}
					JSP	 	& D6:107--108; J1:257--258		& JSR 		& 288 \\
					DC	 	& ---		& HC 		& 3:23 \\
					HDDC 	& ---		& TCHC 		& 3:23 \\
					TPJS 	& 118			& 	 		&  \\
					\end{tabular}
					% HDDC = woodford's DC diss; JSR = Marquardt's JS Revelations; TCHC = Vogel HC
				
				\begin{description}
					\item[Online Access] \href{https://www.josephsmithpapers.org/paper-summary/revelation-17-april-1838/1}{Here}
					% \item[Analysis] \hyperref[XX]{Here}
				\end{description}
				
				\begin{description}
					\item[Background]
				\end{description}
				
					% It might also be useful to give a two sentence context of the period: e.g., "This speech was given in Nauvoo to a small congregation one month prior to the destruction of the Nauvoo Expositor. These and other tensions may have been on the prophet's mind when he gave the speech."
				
				\begin{description}
					\item[Original Source]
				\end{description}
				
					% brief description of original source material, with reference to JSP or somewhere else for more info
					% Background material: it might be helpful to have a note that says whether a given document was presented as a speech, was written in a journal, etc.
					% Also a comment here about how reliable the source might be, or what shape it is in, if there are large omissions or parts that are hard to understand, and so on.
				
				\begin{description}
					\item[Text]
				\end{description}
				
					Revelation given to Brigham Young at Far West April 17\supersu{th} 1838. Verrily thus Saith the Lord, Let my Servant Brigham Young go unto the place which he has baught on Mill Creek and there provide for his family until an effectual door is op[e]ned for the suport of his family untill I shall command [him] to go hence, and not to leave his family untill they are amply provided for Amen.------
					
		% -----
		\subsubsection{Revelation, 26 April 1838 [D\&C 115]}
		% -----
			\label{EMS-1838-JS-26-APR-1838}
			% \label{EMS-1830-JS-DAY-3LETTERMONTH-YEAR}
			
			\begin{description}
					\item[Print Sources]
				\end{description}

					\begin{tabular}{p{0.5in}p{1in}p{0.5in}p{1in}}
					JSP	 	& D6:112--118; J1:258--260				& JSR 		& 289--291 \\
					DC	 	& Section 115		& HC 		& 3:23--25 \\
					HDDC 	& 1506--1516		& TCHC 		& 3:24 \\
					\end{tabular}
					% HDDC = woodford's DC diss; JSR = Marquardt's JS Revelations; TCHC = Vogel HC
				
				\begin{description}
					\item[Online Access] \href{https://www.josephsmithpapers.org/paper-summary/revelation-26-april-1838-dc-115/1}{Here}
					% \item[Analysis] \hyperref[DC115]{Here}
				\end{description}
				
				\begin{description}
					\item[Background]
				\end{description}
				
					% It might also be useful to give a two sentence context of the period: e.g., "This speech was given in Nauvoo to a small congregation one month prior to the destruction of the Nauvoo Expositor. These and other tensions may have been on the prophet's mind when he gave the speech."
				
				\begin{description}
					\item[Original Source]
				\end{description}
				
					% brief description of original source material, with reference to JSP or somewhere else for more info
					% Background material: it might be helpful to have a note that says whether a given document was presented as a speech, was written in a journal, etc.
					% Also a comment here about how reliable the source might be, or what shape it is in, if there are large omissions or parts that are hard to understand, and so on.
				
				\begin{description}
					\item[Text]
				\end{description}
				
					Revelation given in Far West, April 26<\supersu{th}>, 1838, Making known the will of God, concerning the building up of this place and of the Lord's house \&c.
					
					Verrily thus Saith the Lord unto you my Servant Joseph Smith Jr. and also my Servant Sidney Rigdon, and also my Servant Hyrum Smith, and your counselors who are and who shall be hereafter appointed, and also unto my Servant Edward Partridge and his Councilors, and also unto my faithfull Servants, who are of the High Council of my Church in zion (for thus it shall be called) and unto all the Elders and people of my Church of Jesus Christ of Latter Day Saints, Scattered abroad in all the world, For thus shall my Church be called in the Last days even the Church of Jesus Christ of Latter Day Saints, Verrily I say unto you all; arise and shine forth \sout{forth} that thy light may be a standard for the nations and that thy gathering to-gether upon the land of zion and upon her stakes may be for a defence and for a reffuge from the storm and from wrath when it shall be poured out without mixture upon the whole Earth, Let the City Far West, be a holy and consecrated land unto me, and <it shall> be called <most> holy for the ground upon which thou standest is holy Therefore I command you to build an house unto me for the gathering together\sout{ing} of my Saints that they may worship me, and let there be a begining of this work; and a foundation and a preparatory work, this following Summer; and let the begining be made on the 4\supersu{th} day of July next; and from that time forth let my people labour diligently to build an house, unto my name, and in one year from this day, let them recommence laying the foundation of my house; thus let them from that time forth laibour diligently untill it shall be finished, from the Corner Stone thereof unto the top thereof, untill there shall not any thing remain that is not finished.
					
					Verrily I say unto you let not my servant Joseph neither my Servant Sidney, neither my Servant Hyrum, get in debt any more for the building of an house unto my name. But let my house be built unto my name according to the pattern which I will shew unto them, and if my people build it not according to the pattern which I Shall shew unto their presidency, I will not accept it at their hands, But if my people do build it according to the pattern which I shall shew unto their presidency, even my servant Joseph and his Councilors; then I will accept it at the hands of my people, And again; Verrily I say unto you it is my will, that the City Far West should be built up spedily, by the gathering of my Saints, and also that other places should be appointed for stakes in the regions round about as they shall be manifested unto my Servant Joseph from time to time. For behold I will be with him and I will Sanctify him before the people for unto him have I given the Keys of this Kingdom and	ministry even so--- Amen.
					
		% -----
		\subsubsection{Record, 28 April 1838}
		% -----
			\label{EMS-1836-JS-28-APR-1838}
			% \label{EMS-1830-JS-DAY-3LETTERMONTH-YEAR}
			
			\begin{description}
					\item[Print Sources]
				\end{description}

					\begin{tabular}{p{0.5in}p{1in}p{0.5in}p{1in}}
					JSP	 	& J1:261--263		& JSR 		& --- \\
					DC	 	& ---				& HC 		& --- \\
					HDDC 	& ---				& TCHC 		& --- \\
					\end{tabular}
					% HDDC = woodford's DC diss; JSR = Marquardt's JS Revelations; TCHC = Vogel HC
				
				\begin{description}
					\item[Online Access] \href{https://www.josephsmithpapers.org/paper-summary/journal-march-september-1838/20}{Here}
					% \item[Analysis] \hyperref[XX]{Here}
				\end{description}
				
				\begin{description}
					\item[Background]
				\end{description}
				
					% It might also be useful to give a two sentence context of the period: e.g., "This speech was given in Nauvoo to a small congregation one month prior to the destruction of the Nauvoo Expositor. These and other tensions may have been on the prophet's mind when he gave the speech."
					
					Need to explain the trial.
				
				\begin{description}
					\item[Original Source]
				\end{description}
				
					% brief description of original source material, with reference to JSP or somewhere else for more info
					% Background material: it might be helpful to have a note that says whether a given document was presented as a speech, was written in a journal, etc.
					% Also a comment here about how reliable the source might be, or what shape it is in, if there are large omissions or parts that are hard to understand, and so on.
					
					You cannot tell from the source whether JS said these things.
				
				\begin{description}
					\item[Text]
				\end{description}
				
						\hypertarget{EMS-1836-JS-28-APR-1838-a}%
					...After
						\marginnote{a}%
					Justice had ceased to weild \sout{his} <its> sword, Mercy then advanced to rescue its victom, which inspired the heart of President J. Smith J\uline{r}, \& Geo W. Harris who, with profound elequence \sout{<\&>} with <a> deep \& sublime thought, with clemency of feeling, spoke in faivour of \sout{mercy} the defendant, but in length of time, while mercy appeared to be doing her utmost, in contending against justice, the latter at last gained the ascendency, and took full \sout{power over} <possession of> the mind of the speaker, who leveled a voley of darts, which came upon the old man, like <a> huricane\sout{s} upon the mountain tops, which seemingly, was about to sweep the victom entirely out of the reach of mercy, but amidst the clashing of the sword of Justice, mercy still claimed the victom, and saved him still in the church of Jesus Christ of Latter Day Saints,
						\hypertarget{EMS-1836-JS-28-APR-1838-b}%
					and
						\marginnote{b}%
					in this last kingdom Happy is it for those whose sins (like this mans) goes before them to Judgement, that they may yet repent and be saved in the Kingdom of our God. Council desided, that inasmuch as this man, had confessed his sins, and asked for, forgiveness, and promised to mark well the path of his feet, and do, (inasmuch as lay in his power.) what God, Should require\sout{d} at his hand<s>. accordingly, it was decided, that he give up his license as High Priest, and stand as a member in the Church, this in consequence of his being concidered not capable of dignifying that office, \&c Council Adjourned
					
		% -----
		\subsubsection{Editorial, 4 May 1838}
		% -----
			\label{EMS-1838-JS-4-MAY-1838}
			% \label{EMS-1830-JS-DAY-3LETTERMONTH-YEAR}
			
			\begin{description}
					\item[Print Sources]
				\end{description}

					\begin{tabular}{p{0.5in}p{1in}p{0.5in}p{1in}}
					JSP	 	& D6:		& JSR 		& --- \\
					DC	 	& ---		& HC 		& --- \\
					HDDC 	& ---		& TCHC 		& --- \\
					\end{tabular}
					% HDDC = woodford's DC diss; JSR = Marquardt's JS Revelations; TCHC = Vogel HC
				
				\begin{description}
					\item[Online Access] \href{https://www.josephsmithpapers.org/paper-summary/editorial-4-may-1838/1}{Here}
					% \item[Analysis] \hyperref[XX]{Here}
				\end{description}
				
				\begin{description}
					\item[Background]
				\end{description}
				
					% It might also be useful to give a two sentence context of the period: e.g., "This speech was given in Nauvoo to a small congregation one month prior to the destruction of the Nauvoo Expositor. These and other tensions may have been on the prophet's mind when he gave the speech."
				
				\begin{description}
					\item[Original Source]
				\end{description}
				
					% brief description of original source material, with reference to JSP or somewhere else for more info
					% Background material: it might be helpful to have a note that says whether a given document was presented as a speech, was written in a journal, etc.
					% Also a comment here about how reliable the source might be, or what shape it is in, if there are large omissions or parts that are hard to understand, and so on.
				
				\begin{description}
					\item[Text]
				\end{description}
				
						\hypertarget{EMS-1838-JS-4-MAY-1838-a}%
					F\textit{ar West, May}
						\marginnote{a}%
					1838.

					Notwithstanding all the efforts of the enemies to the truth, both from without and within, to the contrary, we are enabled to present this Journal, to the patrons, with the prospect of being able to continue it in time to come, without interruption.

					Great have been the exertions of the opposers to righteousness, to prevent us from sending abroad the doctrines of the church to the world: every effort has been used by the combined influence of all classes of enemies, and of all sects and parties of religion; and of those who are opposed to it, in all its forms to prevent it.

					It is indeed somewhat unexpected to us, to be able to commence printing the Journal again so soon; but the general interest felt in it by the Saints in general, soon, in a degree, repaired the loss which was suffered in the burning of the press in Kirtland; and another establishment, by the exertions of the Saints in Far West, has been obtained, sufficiently large, to print the Journal; and soon will be greatly enlarged, so as to do all the printing necessary, for the whole church.

						\hypertarget{EMS-1838-JS-4-MAY-1838-b}%
					We
						\marginnote{b}%
					have no doubt, but liberal minded men will continue to aid with their means until the establishment will be sufficiently supplied with means to make the largest of the kind, any where in the region of country where it is located.

					In this place, the church is as pleasantly situated as could be expected, taking into consideration their circumstances, as the settlement here is but about eighteen months old, and the first settlers had been driven from their homes, and all their property destroyed, and had to come here without any thing.---But to their honor it may be said, that few people on earth have endured the same degree of persecution, with the same patience.

					Nothing discouraged by the great afflictions and tribulations which they have had to endure for Christ's sake. They united with all their powers, to turn a solitary place into a fruitful field---we do not say a wilderness, for there is not a sufficiency of timber to make it a wilderness---and have exceeded the highest expectations of the most enthusiastic.

					Large bodies of land have been, and are now putting under cultivation.

						\hypertarget{EMS-1838-JS-4-MAY-1838-c}%
					We
						\marginnote{c}%
					might venture an assertion on this point, and that, without the fear of contradiction by those who are acquainted with the settlements in this vicinity, and that is, no part of the world can produce a superior to Caldwell County, if an equal. Eighteen months since without scarcely an inhabitant: at this time the city of ``Far West,'' the county seat, has one hundred and fifty houses, and almost the whole county is taken up, or all that part of it, which can be conveniently settled for want of timber: and large bodies of it are now under cultivation.

					An enconium too high, cannot be placed upon the heads of the enterprising and industrious habits of the people of this county. They are fast making for themselves, and their posterity after them, as beautiful, interesting, and as profitable homes, as can be in any country.

					In a very few years, and it will be said with propriety, ``that the solitary place has become glad for them;'' and we can say, that the people will be as glad for it.

						\hypertarget{EMS-1838-JS-4-MAY-1838-d}%
					This
						\marginnote{d}%
					town ``Far West'' is situated in Caldwell county Missouri, in the midst of a prairie of very rich soil. It is an elevated piece of land, and has a commanding view of the surrounding country for many miles, in every direction. On the north, about one mile passes Shoal Creek, a heavy stream which has many water privileges on it. On the south, a little more than half a mile, runs Goose Creek, a tributary of Shoal. It also is large enough to admit of water-works.

					To all appearance the country is healthy, and the farming interest is equal to that in any part of the world; and the means of living are very easily obtained, not even luxuries excepted.

					From this to the territorial line on the north, is from eighty to one hundred miles, and to the line on the west, twenty five or upwards, or what was the territorial line, before the purchase of what is called the Platt[e] and Nodawa[y] countries, or rather Notawa, which signifies rattle snake.

					It will be seen by this, that this town is situated in the north west corner of the State of Missouri, in the 40th deg. of north latitude. The land is rolling and generally dry; at least, there are no more wet lands, than are necessary for grazing purposes, when the country becomes all subdued.

						\hypertarget{EMS-1838-JS-4-MAY-1838-e}%
					The
						\marginnote{e}%
					Saints here are at perfect peace with all the surrounding inhabitants, and persecution is not so much as once named among them: every man can attend to business without fear or excitement, or being molested in any wise. There are many of the inhabitants of this town, who own lands in the vicinity, and are at this time busily engaged in cultivating them. Hundreds of acres of corn have been planted already, in our immediate neighborhood; and hundreds of acres more are now being planted. (This is the fourth day of May).

					The crops of wheat are very promising, and the prospect is that we will have an abundant harvest. The vast quantities of provision purchased, in this upper country by the United States, for the use of the Garrison, and also for the Indians, have made all kinds of provision dear, and somewhat scarce. Corn is fifty cents per bushel; wheat one dollar; pork from eight, to ten dollars per cwt.; and all kinds of provision on a par with these.

						\hypertarget{EMS-1838-JS-4-MAY-1838-f}%
					Perhaps
						\marginnote{f}%
					it might be thought by some necessary, that we should say something about the affairs of Kirtland.---The burning of the printing office there \&c. But it is now, as in former days. In former days the destroyers of the Saints' property were of the baser sort of mankind, even so it is now. And as the Saints in former days considered a formal notice of them, beneath both their character and standing, so do the Saints in like manner now. Only say as they did; ``That a gang of the baser sort, burned and wasted our property to the utmost of their power'' regardless of law, justice, or humanity, and were upheld in their wickedness, by those who were like the high priest in Paul's day, who though, he sat to judge after the law, commanded Paul to be smitten contrary to law. So it was with our persecutors in the east: for notwithstanding they sat to judge after the law, yet, commanded they our property to be destroyed contrary to law.

					And as Paul and Barnabas did at Iconium. So did we at Kirtland.---``When there was an assault made, both of the Gentiles, and also of the Jews, with their rulers, to use them despitefully, and to stone them, they were ware of it, and fled into Lystria and Derbe, cities of Lyconia, and unto the region that lieth round about. And there they preached the gospel.''

						\hypertarget{EMS-1838-JS-4-MAY-1838-g}%
					So
						\marginnote{g}%
					we did in like manner, taking them for our example. When there was an assault being made, of liars, thieves, and religionists, with their rulers all combined, we were aware of it, and fled to ``Far West,'' and are here preaching the gospel whereunto we are called by the power of God. Let so much suffice for Kirtland.

					We have the gratification of saying to the Elders abroad, that we hope to be able to furnish the Journal regularly, from hence forth, as long as it may be thought wisdom to continue it. And we hope on their part, they will use all their exertions to give it circulation.

					The enemies have made so many attempts to destroy us, and always failed, that we now just laugh at them for fools, as the God of heaven said he would at their calamity.
					
		% -----
		\subsubsection{Discourse, 6 May 1838}
		% -----
			\label{EMS-1838-JS-6-MAY-1838}
			% \label{EMS-1830-JS-DAY-3LETTERMONTH-YEAR}
			
			\begin{description}
					\item[Print Sources]
				\end{description}

					\begin{tabular}{p{0.5in}p{1in}p{0.5in}p{1in}}
					JSP	 	& J1:266; D6:133--134		& JSR 		& --- \\
					DC	 	& ---		& HC 		& 3:27 \\
					HDDC 	& ---		& TCHC 		& 3:28 \\
					TPJS 	& 118--119	& 	 		&  \\ 
					\end{tabular}
					% HDDC = woodford's DC diss; JSR = Marquardt's JS Revelations; TCHC = Vogel HC
				
				\begin{description}
					\item[Online Access] \href{https://www.josephsmithpapers.org/paper-summary/journal-march-september-1838/24}{Here}
					% \item[Analysis] \hyperref[XX]{Here}
				\end{description}
				
				\begin{description}
					\item[Background]
				\end{description}
				
					% It might also be useful to give a two sentence context of the period: e.g., "This speech was given in Nauvoo to a small congregation one month prior to the destruction of the Nauvoo Expositor. These and other tensions may have been on the prophet's mind when he gave the speech."
				
				\begin{description}
					\item[Original Source]
				\end{description}
				
					% brief description of original source material, with reference to JSP or somewhere else for more info
					% Background material: it might be helpful to have a note that says whether a given document was presented as a speech, was written in a journal, etc.
					% Also a comment here about how reliable the source might be, or what shape it is in, if there are large omissions or parts that are hard to understand, and so on.
				
				\begin{description}
					\item[Text]
				\end{description}
				
						\hypertarget{EMS-1838-JS-6-MAY-1838-a}%
					This
						\marginnote{a}%
					day, President Smith. delivered a discourse. to the people. Showing, or setting forth the evils that existed, and would exist, by reason of hasty Judgement or dessisions upon any subject, given by any people. or in judgeing before they hear both sides of the question, He also cautioned them against men men, who should come here whining and grouling about their money, because they had helpt the saints and bore some of the burden with others. and thus thinking that others, (who are still poorer and who have still bore greater burden than themselves) aught to make up their loss \&c. And thus he cautioned them to beware of them for here and there they through [throw?] out foul insinuations, to level as it were a dart to <the> best interests of the Church, \& if possible to destroy the Characters of its Presidency
						\hypertarget{EMS-1838-JS-6-MAY-1838-b}%
					He
						\marginnote{b}%
					also instructed the Church, in the mistories of the Kingdom of God; giving them a history of the Plannets \&c. and of Abrahams writings upon the Plannettary System \&c.  In the after part of the day Pres\supers{t.} Smith spoke upon different subjects he dwelt some upon the Subject of Wisdom, \& upon the word of Wisdom. \&c.
					
		% -----
		\subsubsection{Questions and Answers, 8 May 1838}
		% -----
			\label{EMS-1838-JS-8-MAY-1838}
			% \label{EMS-1830-JS-DAY-3LETTERMONTH-YEAR}
			
			\begin{description}
					\item[Print Sources]
				\end{description}

					\begin{tabular}{p{0.5in}p{1in}p{0.5in}p{1in}}
					JSP	 	& D6:139--145		& JSR 		& --- \\
					DC	 	& ---		& HC 		& 3:28--30 \\
					HDDC 	& ---		& TCHC 		& 3:29--31 \\
					TPJS 	& 119--121			& 	 		&  \\
					\end{tabular}
					% HDDC = woodford's DC diss; JSR = Marquardt's JS Revelations; TCHC = Vogel HC
				
				\begin{description}
					\item[Online Access] \href{https://www.josephsmithpapers.org/paper-summary/questions-and-answers-8-may-1838/1}{Here}
					% \item[Analysis] \hyperref[XX]{Here}
				\end{description}
				
				\begin{description}
					\item[Background]
				\end{description}
				
					% It might also be useful to give a two sentence context of the period: e.g., "This speech was given in Nauvoo to a small congregation one month prior to the destruction of the Nauvoo Expositor. These and other tensions may have been on the prophet's mind when he gave the speech."
				
				\begin{description}
					\item[Original Source]
				\end{description}
				
					% brief description of original source material, with reference to JSP or somewhere else for more info
					% Background material: it might be helpful to have a note that says whether a given document was presented as a speech, was written in a journal, etc.
					% Also a comment here about how reliable the source might be, or what shape it is in, if there are large omissions or parts that are hard to understand, and so on.
				
				\begin{description}
					\item[Text]
				\end{description}
				
						\hypertarget{EMS-1838-JS-8-MAY-1838-a}%
					In
						\marginnote{a}%
					obedience to our promise, we give the following answers to questions, which were asked in the last number of the Journal.

					Question 1st. Do you believe the b[i]ble?

					Answer. If we do, we are the only people under heaven that does. For there are none of the religious sects of the day that do.

					Question 2nd. Wherein do you differ from other sects?

					Answer. Because we believe the bible, and all other sects profess to believe their interpretations of the bible, and their creeds.

					Question 3rd. Will every body be damned but Mormons?

					Answer. Yes, and a great portion of them, unless they repent and work righteousness.

					Question 4th. How, and where did you obtain the book of Mormon?

					Answer. Moroni, the person who deposited the plates, from whence the book of Mormon was translated, in a hill in Manchester, Ontario County New York, being dead, and raised again therefrom, appeared unto me, and told me where they were; and gave me directions how to obtain them. I obtained them, and the Urim and Thummim with them; by the means of which, I translated the plates; and thus came the book of Mormon.

						\hypertarget{EMS-1838-JS-8-MAY-1838-b}%
					Question
						\marginnote{b}%
					5th. Do you believe Joseph Smith Jr. to be a prophet?

					Answer. Yes, and every other man who has the testimony of Jesus. ``For the testimony of Jesus, is the spirit of prophecy.''--- Rev. 19:10.

					Question 6th. Do the Mormons believe in having all things common?

					Answer. No.

					Question 7th. Do the Mormons believe in having more wives than one.

					Answer. No, not at the same time. But they believe, that if their companion dies, they have a right to marry again. But we do disapprove of the custom which has gained in the world, and has been practised among us, to our great mortification, of marrying in five or six weeks, or even in two or three months after the death of their companion.

					We believe that due respect ought to be had, to the memory of the dead, and the feelings of both friends and children.

					Question 8th. Can they raise the dead.

					Answer. No, nor any other people that now lives or ever did live. But God can raise the dead through man, as an instrument.

					Question 9th. What signs do Jo Smith give of his divine mission.

					Answer. The signs which God is pleased to let him give: according as his wisdom thinks best: in order that he may judge the world agreably to his own plan.

						\hypertarget{EMS-1838-JS-8-MAY-1838-c}%
					Question
						\marginnote{c}%
					10. Was not Jo Smith a money digger.

					Answer. Yes, but it was never a very profitable job to him, as he only got fourteen dollars a month for it.

					Question 11th. Did not Jo Smith steal his wife.

					Answer. Ask her; she was of age, she can answer for herself.

					Question 12th. Do the people have to give up their money, when they join his church.

					Answer. No other requirement than to bear their proportion of the expenses of the church, and support the poor.

					Question 13th. Are the Mormons abolitionists.

					Answer. No, unless delivering the people from priest-craft, and the priests from the prower of satan, should be considered such.--- But we do not believe in setting the Negroes free.

					Question 14th. Do they not stir up the Indians to war and to commit depredations.

					Answer. No, and those who reported the story, knew it was false when they put it into circulation. These and similar reports, are pawned upon the people by the priests, and this is the reason why we ever thought of answering them.

					Question 15th. Do the Mormons baptize in the name of Jo Smith.

					Answer. No, but if they did, it would be as valid as the baptism administered by the sectarian priests.

						\hypertarget{EMS-1838-JS-8-MAY-1838-d}%
					Question
						\marginnote{d}%
					16th. If the Mormon doctrine is true what has become of all those who have died since the days of the apostles.

					Answer. All those who have not had an opportunity of hearing the gospel, and being administered to by an inspired man in the flesh, must have it hereafter, before they can be finally judged.

					Question 17th. Does not Jo Smith profess to be Jesus Christ.

					Answer. No, but he professes to be his brother, as all other saints have done, and now do.--- Matthew, 12:49, 50--- And he stretched forth his hand toward his disciples and said, Behold my mother and my brethren: For whosoever shall do the will of my father which is in heaven, the same is my brother, and sister, and mother.

					Question 18th. Is there any thing in the Bible which lisences you to believe in revelation now a days.

					Answer. Is there any thing that does not authorize us to believe so; if there is, we have, as yet, not been able to find it.
					
						\hypertarget{EMS-1838-JS-8-MAY-1838-e}%
					Question
						\marginnote{e}%
					19th. Is not the cannon of the Scriptures full.
					
					Answer. If it is, there is a great defect in the book, or else it would have said so.

					Question 20th. What are the fundamental principles of your religion.

					Answer. The fundamental principles of our religion is the testimony of the apostles and prophets concerning Jesus Christ, ``that he died, was buried, and rose again the third day, and ascended up into heaven;'' and all other things are only appendages to these, which pertain to our religion.

					But in connection with these, we believe in the gift of the Holy Ghost, the power of faith, the enjoyment of the spiritual gifts according to the will of God, the restoration of the house of Israel, and the final triumph of truth.
					
		% -----
		\subsubsection{Record, 9 May 1838}
		% -----
			\label{EMS-1838-JS-9-MAY-1838}
			% \label{EMS-1830-JS-DAY-3LETTERMONTH-YEAR}
			
			\begin{description}
					\item[Print Sources]
				\end{description}

					\begin{tabular}{p{0.5in}p{1in}p{0.5in}p{1in}}
					JSP	 	& J1:267		& JSR 		& --- \\
					DC	 	& ---				& HC 		& --- \\
					HDDC 	& ---				& TCHC 		& --- \\
					\end{tabular}
					% HDDC = woodford's DC diss; JSR = Marquardt's JS Revelations; TCHC = Vogel HC
				
				\begin{description}
					\item[Online Access] \href{https://www.josephsmithpapers.org/paper-summary/journal-march-september-1838/25}{Here}
					% \item[Analysis] \hyperref[XX]{Here}
				\end{description}
				
				\begin{description}
					\item[Background]
				\end{description}
				
					% It might also be useful to give a two sentence context of the period: e.g., "This speech was given in Nauvoo to a small congregation one month prior to the destruction of the Nauvoo Expositor. These and other tensions may have been on the prophet's mind when he gave the speech."
				
				\begin{description}
					\item[Original Source]
				\end{description}
				
					% brief description of original source material, with reference to JSP or somewhere else for more info
					% Background material: it might be helpful to have a note that says whether a given document was presented as a speech, was written in a journal, etc.
					% Also a comment here about how reliable the source might be, or what shape it is in, if there are large omissions or parts that are hard to understand, and so on.
				
				\begin{description}
					\item[Text]
				\end{description}
				
					This day, the Presidency, attended the funeral of James Marsh, And Pres\supersu{t}\supers{.} Smith was requested to preach the funeral discourse, and accordingly complied, and we were greatly edified upon the occasion,
					
		% -----
		\subsubsection{Record, 11 May 1838}
		% -----
			\label{EMS-1838-JS-11-MAY-1838}
			% \label{EMS-1830-JS-DAY-3LETTERMONTH-YEAR}
			
			\begin{description}
					\item[Print Sources]
				\end{description}

					\begin{tabular}{p{0.5in}p{1in}p{0.5in}p{1in}}
					JSP	 	& J1:268		& JSR 		& --- \\
					DC	 	& ---				& HC 		& --- \\
					HDDC 	& ---				& TCHC 		& --- \\
					\end{tabular}
					% HDDC = woodford's DC diss; JSR = Marquardt's JS Revelations; TCHC = Vogel HC
				
				\begin{description}
					\item[Online Access] \href{https://www.josephsmithpapers.org/paper-summary/journal-march-september-1838/26}{Here}
					% \item[Analysis] \hyperref[XX]{Here}
				\end{description}
				
				\begin{description}
					\item[Background]
				\end{description}
				
					% It might also be useful to give a two sentence context of the period: e.g., "This speech was given in Nauvoo to a small congregation one month prior to the destruction of the Nauvoo Expositor. These and other tensions may have been on the prophet's mind when he gave the speech."
				
				\begin{description}
					\item[Original Source]
				\end{description}
				
					% brief description of original source material, with reference to JSP or somewhere else for more info
					% Background material: it might be helpful to have a note that says whether a given document was presented as a speech, was written in a journal, etc.
					% Also a comment here about how reliable the source might be, or what shape it is in, if there are large omissions or parts that are hard to understand, and so on.
				
				\begin{description}
					\item[Text]
				\end{description}
				
					W\supersu{m}\supers{.} E. Mc.Lellin, also said the same. He further said he had no confidence in the heads of the Church, beleiving they had transgressed, and got out of the way, and consequently <he> left of[f] praying and keeping the commandments of God, and went his own way, and indulged himself in his lustfull desires. But when he heard, that the first presidency, had made a general settlement and acknowleged their sins, he then began to pray again, and to keep the commandments of God. Though when interogated by Pres\supers{t} smith he said he had seen nothing out of the way himself but it was heresay, and thus he judged from heresay. But we are constrained to say, O!! foolish Man! what excuse is that thou renderest, for thy sins, that because thou hast heard of some mans transgression, that thou shouldest leave thy God, and forsake thy prayers, and turn to those things that thou knowest to be contrary to the will of God, we say unto thee, and to all such, beware! beware! for God will bring the[e] into Judgement for thy sins.
					
		% -----
		\subsubsection{Record, 19 May 1838}
		% -----
			\label{EMS-1838-JS-19-MAY-1838}
			% \label{EMS-1830-JS-DAY-3LETTERMONTH-YEAR}
			
			\begin{description}
					\item[Print Sources]
				\end{description}

					\begin{tabular}{p{0.5in}p{1in}p{0.5in}p{1in}}
					JSP	 	& ---				& JSR 		& --- \\
					DC	 	& ---				& HC 		& --- \\
					HDDC 	& ---				& TCHC 		& --- \\
					TPJS 	& 122				&  &  \\
					\end{tabular}
					% HDDC = woodford's DC diss; JSR = Marquardt's JS Revelations; TCHC = Vogel HC
				
				\begin{description}
					\item[Online Access] \href{https://www.josephsmithpapers.org/paper-summary/journal-march-september-1838/29}{Here}
					% \item[Analysis] \hyperref[XX]{Here}
				\end{description}
				
				\begin{description}
					\item[Background]
				\end{description}
				
					% It might also be useful to give a two sentence context of the period: e.g., "This speech was given in Nauvoo to a small congregation one month prior to the destruction of the Nauvoo Expositor. These and other tensions may have been on the prophet's mind when he gave the speech."
				
				\begin{description}
					\item[Original Source]
				\end{description}
				
					% brief description of original source material, with reference to JSP or somewhere else for more info
					% Background material: it might be helpful to have a note that says whether a given document was presented as a speech, was written in a journal, etc.
					% Also a comment here about how reliable the source might be, or what shape it is in, if there are large omissions or parts that are hard to understand, and so on.
				
				\begin{description}
					\item[Text]
				\end{description}
						
						\hypertarget{EMS-1838-JS-19-MAY-1838-a}%
					The
						\marginnote{a}%
					next morning we struck our tents, and marched crossed Grand river at the mouth of Honey Creek at a place called Nelsons ferry, Grand River is a\sout{s} large beautifull deep and rapid stream and will undoubtedly admit of steam Boat and other water craft navigation, and at the mouth of honey creek is a splendid harbour for the safety of such crafts, and also for landing freight We next kept up the river mostly in the timber for ten miles, untill we came to Col. Lyman Wight's who lives at the foot of Tower Hill, a name appropriated by Pres\supers{t} smith, in consequence of the remains of an old Nephitish Alter an Tower, \sout{In the after} where we camped for the sabath,
						\hypertarget{EMS-1838-JS-19-MAY-1838-b}%
					In
						\marginnote{b}%
					the after part of the day, Prest<\supers{s.}> smith and Rigdon and myself, went to Wights. Ferry about a half mile from this place up the river, for the purpose of selecting and laying claims to city plott near said Ferry, in Davis [Daviess] County Township 60, Range 27 \& 28, and Sections 25, 36, 31, 30, which was called Spring Hill a name appropriated by the bretheren present, But after wards named by the mouth of [the] Lord and was called Adam Ondi Awmen [Adam-ondi-Ahman], because said he it is the place where Adam shall come to visit his people, or the Ancient of days shall sit as spoken of by Daniel the Prophet,
					
		% -----
		\subsubsection{Minutes, 28 June 1838}
		% -----
			\label{EMS-1838-JS-28-JUN-1838}
			% \label{EMS-1830-JS-DAY-3LETTERMONTH-YEAR}
			
			\begin{description}
					\item[Print Sources]
				\end{description}

					\begin{tabular}{p{0.5in}p{1in}p{0.5in}p{1in}}
					JSP	 	& D6:162--167		& JSR 		& --- \\
					DC	 	& ---		& HC 		& 3:38--39 \\
					HDDC 	& ---		& TCHC 		& 3:38--39 \\
					\end{tabular}
					% HDDC = woodford's DC diss; JSR = Marquardt's JS Revelations; TCHC = Vogel HC
				
				\begin{description}
					\item[Online Access] \href{https://www.josephsmithpapers.org/paper-summary/minutes-28-june-1838/1}{Here}
					% \item[Analysis] \hyperref[XX]{Here}
				\end{description}
				
				\begin{description}
					\item[Background]
				\end{description}
				
					% It might also be useful to give a two sentence context of the period: e.g., "This speech was given in Nauvoo to a small congregation one month prior to the destruction of the Nauvoo Expositor. These and other tensions may have been on the prophet's mind when he gave the speech."
				
				\begin{description}
					\item[Original Source]
				\end{description}
				
					% brief description of original source material, with reference to JSP or somewhere else for more info
					% Background material: it might be helpful to have a note that says whether a given document was presented as a speech, was written in a journal, etc.
					% Also a comment here about how reliable the source might be, or what shape it is in, if there are large omissions or parts that are hard to understand, and so on.
				
				\begin{description}
					\item[Text]
				\end{description}
				
					...A conference meeting of Elders, and members, of the church of Christ of Latter Day Saints, was held in this place, this day, for the purpose of organziing this stake of Zion, called Adam-ondi-ahman. The meeting convened at 10 o'clock A. M. in the grove near the house of elder Lyman Wight. President Joseph Smith Jr. was called. to the chair, who explained the object of the meeting, which was to organize a Presidency, and High Council, to preside over this stake of Zion, and attend to the affairs of the church in Daviess county...
					
					...After the ordination of the counsellors, who had not previously been ordained to the high priesthood. President J. Smith Jr. made remarks by way of charge to the Presidents and counsellors, instructing them in the duty of their callings, and the responsibility of their stations; exhorting them to be cautious and deliberate, in all their councils, and to be careful to act in righteousness in all things...
					
		% -----
		\subsubsection{Revelation, 8 July 1838--A [D\&C 118]}
		% -----
			\label{EMS-1838-JS-8-JUL-1838-A}
			% \label{EMS-1830-JS-DAY-3LETTERMONTH-YEAR}
			
			\begin{description}
					\item[Print Sources]
				\end{description}

					\begin{tabular}{p{0.5in}p{1in}p{0.5in}p{1in}}
					JSP	 	& D6:175--180; J1:284--285				& JSR 		& 292 \\
					DC	 	& Section 118		& HC 		& 3:46--47 \\
					HDDC 	& 1536--1550		& TCHC 		& 3:46--47 \\
					\end{tabular}
					% HDDC = woodford's DC diss; JSR = Marquardt's JS Revelations; TCHC = Vogel HC
				
				\begin{description}
					\item[Online Access] \href{https://www.josephsmithpapers.org/paper-summary/revelation-8-july-1838-a-dc-118/1}{Here}
					% \item[Analysis] \hyperref[DC118]{Here}
				\end{description}
				
				\begin{description}
					\item[Background]
				\end{description}
				
					% It might also be useful to give a two sentence context of the period: e.g., "This speech was given in Nauvoo to a small congregation one month prior to the destruction of the Nauvoo Expositor. These and other tensions may have been on the prophet's mind when he gave the speech."
				
				\begin{description}
					\item[Original Source]
				\end{description}
				
					% brief description of original source material, with reference to JSP or somewhere else for more info
					% Background material: it might be helpful to have a note that says whether a given document was presented as a speech, was written in a journal, etc.
					% Also a comment here about how reliable the source might be, or what shape it is in, if there are large omissions or parts that are hard to understand, and so on.
				
				\begin{description}
					\item[Text]
				\end{description}
				
					Revelation given July 8\supersu{th}\supers{.} 1838 at Far West Caldwell Co Mo
					
					Shew unto us thy will O Lord concerning the Twelve---
					
					Ans--- Verily thus saith the Lord let a conference be held immediately let the Twelve be organized and let men be appointed to supply the place of those who are fallen let my servant Thomas remain for a \sout{little} season in the Land of Zion to publish my word let the remainder continue to prech from that hour and if they will do this in all lowliness of heart in meekness and pureness and long sufferring, I the Lord God give unto them a promise that I will provide for their families and an effectual door shall be opened for \sout{their families} them from henceforth and next spring let them depart to go over the great waters and there promulge my gospel the fulness thereof and to bear record of my name let them take leave of my saints in the City Far West on the 26\supersu{th}\supers{.} day of April next on the building spot of my house saith the Lord let my servant John Taylor and also my servant John E. Page and also my servant Wilford Woodruff and also my servant Willard Richards be appointed to fill the places of those who have fallen and be officially notified of their appointment [\textit{11 lines blank}]
					
		% -----
		\subsubsection{Revelation, 8 July 1838--B}
		% -----
			\label{EMS-1838-JS-8-JUL-1838-B}
			% \label{EMS-1830-JS-DAY-3LETTERMONTH-YEAR}
			
			\begin{description}
					\item[Print Sources]
				\end{description}

					\begin{tabular}{p{0.5in}p{1in}p{0.5in}p{1in}}
					JSP	 	& D6:181--183; J1:285--288		& JSR 		& 292--293 \\
					DC	 	& ---		& HC 		& 3:44 \\
					HDDC 	& ---		& TCHC 		& 3:44 \\
					\end{tabular}
					% HDDC = woodford's DC diss; JSR = Marquardt's JS Revelations; TCHC = Vogel HC
				
				\begin{description}
					\item[Online Access] \href{https://www.josephsmithpapers.org/paper-summary/revelation-8-july-1838-b/1}{Here}
					% \item[Analysis] \hyperref[XX]{Here}
				\end{description}
				
				\begin{description}
					\item[Background]
				\end{description}
				
					% It might also be useful to give a two sentence context of the period: e.g., "This speech was given in Nauvoo to a small congregation one month prior to the destruction of the Nauvoo Expositor. These and other tensions may have been on the prophet's mind when he gave the speech."
				
				\begin{description}
					\item[Original Source]
				\end{description}
				
					% brief description of original source material, with reference to JSP or somewhere else for more info
					% Background material: it might be helpful to have a note that says whether a given document was presented as a speech, was written in a journal, etc.
					% Also a comment here about how reliable the source might be, or what shape it is in, if there are large omissions or parts that are hard to understand, and so on.
				
				\begin{description}
					\item[Text]
				\end{description}
				
					Revelation Given July 8. 1838 in Far West M\supers{o.}
					
					O Lord what is thy will concerning W[illiam] W Phelps \& F[rederick] G. Williams
					
					Verily thus saith the Lord in consequence of their transgressions their former standing has been taken away from them, and now, if they will be saved, let them be ordained as Elders in my \sout{vineyard} Church to Preach my gospel, and to travel abroad from land to land \& from place to place to gather mine Elect unto me saith the Lord, and let this be their labours from henceforth \uline{Amen} [\textit{3/4 page blank}]
					
		% -----
		\subsubsection{Revelation, 8 July 1838--C [D\&C 119]}
		% -----
			\label{EMS-1838-JS-8-JUL-1838-C}
			% \label{EMS-1830-JS-DAY-3LETTERMONTH-YEAR}
			
			\begin{description}
					\item[Print Sources]
				\end{description}

					\begin{tabular}{p{0.5in}p{1in}p{0.5in}p{1in}}
					JSP	 	& D6:183--189; J1:288				& JSR 		& 293 \\
					DC	 	& Section 119		& HC 		& 3:44 \\
					HDDC 	& 1551--1558		& TCHC 		& 3:45 \\
					\end{tabular}
					% HDDC = woodford's DC diss; JSR = Marquardt's JS Revelations; TCHC = Vogel HC
				
				\begin{description}
					\item[Online Access] \href{https://www.josephsmithpapers.org/paper-summary/revelation-8-july-1838-c-dc-119/1}{Here}
					% \item[Analysis] \hyperref[DC119]{Here}
				\end{description}
				
				\begin{description}
					\item[Background]
				\end{description}
				
					% It might also be useful to give a two sentence context of the period: e.g., "This speech was given in Nauvoo to a small congregation one month prior to the destruction of the Nauvoo Expositor. These and other tensions may have been on the prophet's mind when he gave the speech."
				
				\begin{description}
					\item[Original Source]
				\end{description}
				
					% brief description of original source material, with reference to JSP or somewhere else for more info
					% Background material: it might be helpful to have a note that says whether a given document was presented as a speech, was written in a journal, etc.
					% Also a comment here about how reliable the source might be, or what shape it is in, if there are large omissions or parts that are hard to understand, and so on.
				
				\begin{description}
					\item[Text]
				\end{description}
				
					Far West July 8\supers{th} 1838

					A revelation

					Question O Lord show unto thy servants how much thou requirest of the properties of thy people for a tithing?

					Answer. Verily thus saith the Lord, I require all their surplus property, to be put into the hands of the bishop of my church of Zion for the building of mine house and for the laying the foundation of Zion and for the priesthood and for the debts of the presidency of my church and this shall be the beginning of the tithing of my people and after that those who have \sout{been} thus been tithed shall pay one tenth of all their interest annually and this shall be a standing law unto them forever for my holy priesthood saith the Lord. Verily I say unto you it shall come to pass that all those who gather unto the land of Zion shall be tithed of their surplus properties and shall observe this law or they shall not be found worthy to abide among you and behold I say unto you if my people observe not this law to keep it holy and by this law sanctify the land of Zion unto me that my statutes and <my> judgements may be kept thereon that it may be most holy, behold verily I say unto you it shall not be a land of Zion unto you[.] And this shall be an ensample unto all the stakes of Zion even so amen

					[\textit{page [2] blank}]
					
		% -----
		\subsubsection{Revelation, 8 July 1838--D [D\&C 120]}
		% -----
			\label{EMS-1838-JS-8-JUL-1838-D}
			% \label{EMS-1830-JS-DAY-3LETTERMONTH-YEAR}
			
			\begin{description}
					\item[Print Sources]
				\end{description}

					\begin{tabular}{p{0.5in}p{1in}p{0.5in}p{1in}}
					JSP	 	& D6:189--190; J1:289				& JSR 		& 293--294 \\
					DC	 	& Section 120		& HC 		& 3:44 \\
					HDDC 	& 1559--1565		& TCHC 		& 3:45 \\
					\end{tabular}
					% HDDC = woodford's DC diss; JSR = Marquardt's JS Revelations; TCHC = Vogel HC
				
				\begin{description}
					\item[Online Access] \href{https://www.josephsmithpapers.org/paper-summary/revelation-8-july-1838-d-dc-120/1}{Here}
					% \item[Analysis] \hyperref[DC120]{Here}
				\end{description}
				
				\begin{description}
					\item[Background]
				\end{description}
				
					% It might also be useful to give a two sentence context of the period: e.g., "This speech was given in Nauvoo to a small congregation one month prior to the destruction of the Nauvoo Expositor. These and other tensions may have been on the prophet's mind when he gave the speech."
				
				\begin{description}
					\item[Original Source]
				\end{description}
				
					% brief description of original source material, with reference to JSP or somewhere else for more info
					% Background material: it might be helpful to have a note that says whether a given document was presented as a speech, was written in a journal, etc.
					% Also a comment here about how reliable the source might be, or what shape it is in, if there are large omissions or parts that are hard to understand, and so on.
				
				\begin{description}
					\item[Text]
				\end{description}
				
					Revelation Given the same day July 8\supers{th} 1838
					
					Making known the disposition of the properties tithed, as named in the preceeding revelation--- \\
					
					Verrily thus saith the Lord, the time has now come that it shall be disposed of, by a council composed of the first Presidency of my Church and of the Bishop and his council and by <my> high Council, and <by> mine own voice unto them saith the Lord, even so Amen.
		
		% -----
		\subsubsection{Revelation, 8 July 1838--E [D\&C 117]}
		% -----
			\label{EMS-1838-JS-8-JUL-1838-E}
			% \label{EMS-1830-JS-DAY-3LETTERMONTH-YEAR}
			
			\begin{description}
					\item[Print Sources]
				\end{description}

					\begin{tabular}{p{0.5in}p{1in}p{0.5in}p{1in}}
					JSP	 	& D6:191--194; J1:289--290				& JSR 		& 294--295 \\
					DC	 	& Section 117		& HC 		& 3:45--46 \\
					HDDC 	& 1523--1535		& TCHC 		& 3:45--46 \\
					\end{tabular}
					% HDDC = woodford's DC diss; JSR = Marquardt's JS Revelations; TCHC = Vogel HC
				
				\begin{description}
					\item[Online Access] \href{https://www.josephsmithpapers.org/paper-summary/revelation-8-july-1838-e-dc-117/1}{Here}
					% \item[Analysis] \hyperref[DC117]{Here}
				\end{description}
				
				\begin{description}
					\item[Background]
				\end{description}
				
					% It might also be useful to give a two sentence context of the period: e.g., "This speech was given in Nauvoo to a small congregation one month prior to the destruction of the Nauvoo Expositor. These and other tensions may have been on the prophet's mind when he gave the speech."
				
				\begin{description}
					\item[Original Source]
				\end{description}
				
					% brief description of original source material, with reference to JSP or somewhere else for more info
					% Background material: it might be helpful to have a note that says whether a given document was presented as a speech, was written in a journal, etc.
					% Also a comment here about how reliable the source might be, or what shape it is in, if there are large omissions or parts that are hard to understand, and so on.
					
					% Note the 8 July 1838 letter to William Marks, which includes a version of this revelation.
				
				\begin{description}
					\item[Text]
				\end{description}
				
					A revelation given at Far West July 8\supers{th} A.D. 1838

					verily thus saith the Lord unto my servants William Marks and N[ewel] K. Whitney, Let them settle up their businss speedily and journey from the land of Kirtland, before I the Lord sendeth snow again upon the ground. Let them awake and arise and come forth and not tarry, for I the Lord commandeth it.--- therefore if they tarry, it shall not be well with them.--- Let them repent of all their sins, and all their covetous desires before me saith the Lord: \sout{And whatsoever remaineth, let it remain in your hands} for what is property unto me saith the Lord. Let the properties at Kirtland be turned out for debts saith the Lord. Let them go saith the Lord; and whatsoever remaineth, let it remain in your hands saith the Lord: for have I not the fowls of Heaven, and also the fish of the sea, and the beasts of the mountains. Have I not made the earth? Do I not hold the destinies of all the armies of the nations of the earth? Therefore will I not make the solitary places to bud and to blossom and to bring forth in abundance saith the Lord. Is there not room enough upon the mountains of Adamondi awman, or upon the plains of Olea Shinihah, or the land where Adam dwelt, that you should not covet that which is but \sout{but} the drop, and neglect the more w[e]ighty matters.--- Therefore come up hither unto the land of my people even Zion. Let my servant William Marks be faithful over a few things, and he shall be ruler over many things. Let him preside in the midst of my people in the city Far West, and let him be blessed with the blessings of my people.--- Let my servant N. K. Whitney be ashamed of the Nicolitans, and of all their secret abominations, and of all his littleness of soul before me saith the Lord, and come up unto the land of Adam ondi awman and be a \sout{man} bishop unto my people saith the Lord, not in name but in deed saith the Lord. And again verily I say unto you I remember my servant Oliver Granger. Behold verily I say unto him, that his name shall be had in sacred rememberance from generation to generation forever and ever saith the Lord. Therefore let him contend earnestly for the redemption of the first presidency of my church saith the Lord; and when he falls he shall rise again; for his sacrifice shall be more sacred to me, than his increase saith the Lord; therefore let him come up hither speedily unto the land of Zion, and in due time he shall be made a merchant unto my name saith the Lord for the benefit of my people: therefore let no man despise my servant Oliver Granger, but let the blessings of my people be upon him forever and ever. Amen.

					And again I say unto you let all the saints in the land of Kirtland remember the Lord their God, and mine house to preserve it holy, and to overthrow the money changers in mine own due time said the Lord [\textit{3/4 page blank}]
					
		% -----
		\subsubsection{Discourse, 29 July 1838}
		% -----
			\label{EMS-1838-JS-29-JUL-1838}
			% \label{EMS-1830-JS-DAY-3LETTERMONTH-YEAR}
			
			\begin{description}
					\item[Print Sources]
				\end{description}

					\begin{tabular}{p{0.5in}p{1in}p{0.5in}p{1in}}
					JSP	 	& D6:209--212		& JSR 		& --- \\
					DC	 	& ---		& HC 		& --- \\
					HDDC 	& ---		& TCHC 		& --- \\
					\end{tabular}
					% HDDC = woodford's DC diss; JSR = Marquardt's JS Revelations; TCHC = Vogel HC
				
				\begin{description}
					\item[Online Access] \href{https://www.josephsmithpapers.org/paper-summary/discourse-29-july-1838/1}{Here}
					% \item[Analysis] \hyperref[XX]{Here}
				\end{description}
				
				\begin{description}
					\item[Background]
				\end{description}
				
					% It might also be useful to give a two sentence context of the period: e.g., "This speech was given in Nauvoo to a small congregation one month prior to the destruction of the Nauvoo Expositor. These and other tensions may have been on the prophet's mind when he gave the speech."
				
				\begin{description}
					\item[Original Source]
				\end{description}
				
					% brief description of original source material, with reference to JSP or somewhere else for more info
					% Background material: it might be helpful to have a note that says whether a given document was presented as a speech, was written in a journal, etc.
					% Also a comment here about how reliable the source might be, or what shape it is in, if there are large omissions or parts that are hard to understand, and so on.
				
				\begin{description}
					\item[Text]
				\end{description}
				
					Brother Joseph Smith then preached, and took his text in 1 Tessalonians, 5:15 to 23d verse. He preached on prophecy, and said that the Spirit of God had appeared to him, with wonderful \textit{light} and \textit{mystery}---in such a manner that we would not believe him, were he to tell us what he had seen; and that he could not say what God did these things for. ``But,'' said he, ``I cannot help it. I know that all the world is threatening my life; but I regard it not, for I am willing to die at any time when God calls for me. I have been beaten, abused, stoned, persecuted, and have had to escape by day and by night. I have been sued at law, and have always proved myself innocent. I have had \textit{twenty-one law-suits}. I am of age; and care not how long I live. Not my will be done, but thine, O Lord!''
					
		% -----
		\subsubsection{Record, 6 August 1838}
		% -----
			\label{EMS-1836-JS-6-AUG-1838}
			% \label{EMS-1830-JS-DAY-3LETTERMONTH-YEAR}
			
			\begin{description}
					\item[Print Sources]
				\end{description}

					\begin{tabular}{p{0.5in}p{1in}p{0.5in}p{1in}}
					JSP	 	& J1:298		& JSR 		& --- \\
					DC	 	& ---				& HC 		& --- \\
					HDDC 	& ---				& TCHC 		& --- \\
					\end{tabular}
					% HDDC = woodford's DC diss; JSR = Marquardt's JS Revelations; TCHC = Vogel HC
				
				\begin{description}
					\item[Online Access] \href{https://www.josephsmithpapers.org/paper-summary/journal-march-september-1838/51}{Here}
					% \item[Analysis] \hyperref[XX]{Here}
				\end{description}
				
				\begin{description}
					\item[Background]
				\end{description}
				
					% It might also be useful to give a two sentence context of the period: e.g., "This speech was given in Nauvoo to a small congregation one month prior to the destruction of the Nauvoo Expositor. These and other tensions may have been on the prophet's mind when he gave the speech."
				
				\begin{description}
					\item[Original Source]
				\end{description}
				
					% brief description of original source material, with reference to JSP or somewhere else for more info
					% Background material: it might be helpful to have a note that says whether a given document was presented as a speech, was written in a journal, etc.
					% Also a comment here about how reliable the source might be, or what shape it is in, if there are large omissions or parts that are hard to understand, and so on.
				
				\begin{description}
					\item[Text]
				\end{description}
				
					Pres\supers{t.} Smith spoke upon the same subject of mooving into Cities to live, according to the order of God, he spoke quite lengthy and then Pr\supers{est.} H[yrum] Smith spoke and endeavoured to impress it upon the \sout{same} minds of the Saints,
					
		% -----
		\subsubsection{Discourse, 12 August 1838}
		% -----
			\label{EMS-1838-JS-12-JUL-1838}
			% \label{EMS-1830-JS-DAY-3LETTERMONTH-YEAR}
			
			\begin{description}
					\item[Print Sources]
				\end{description}

					\begin{tabular}{p{0.5in}p{1in}p{0.5in}p{1in}}
					JSP	 	& D6:213--215		& JSR 		& --- \\
					DC	 	& ---		& HC 		& 3:62--63 \\
					HDDC 	& ---		& TCHC 		& 3:66 \\
					\end{tabular}
					% HDDC = woodford's DC diss; JSR = Marquardt's JS Revelations; TCHC = Vogel HC
				
				\begin{description}
					\item[Online Access] \href{https://www.josephsmithpapers.org/paper-summary/discourse-12-august-1838/1}{Here}
					% \item[Analysis] \hyperref[XX]{Here}
				\end{description}
				
				\begin{description}
					\item[Background]
				\end{description}
				
					% It might also be useful to give a two sentence context of the period: e.g., "This speech was given in Nauvoo to a small congregation one month prior to the destruction of the Nauvoo Expositor. These and other tensions may have been on the prophet's mind when he gave the speech."
				
				\begin{description}
					\item[Original Source]
				\end{description}
				
					% brief description of original source material, with reference to JSP or somewhere else for more info
					% Background material: it might be helpful to have a note that says whether a given document was presented as a speech, was written in a journal, etc.
					% Also a comment here about how reliable the source might be, or what shape it is in, if there are large omissions or parts that are hard to understand, and so on.
				
				\begin{description}
					\item[Text]
				\end{description}
				
					The Prophet exhorted the brethren to be of good cheer, and not to be ``scared at trifles,'' ``that he would tell them when danger was near---that at present he did not even anticipate such a thing---for us to hold ourselves in readiness at a moment's warning, well armed and equipped, to go down to Carroll county and rescue brother [George M.] Hinkle, who is in a suffering condition; and that they must expect to endure trials there.''
					
		% -----
		\subsubsection{Letter to Stephen Post, 17 September 1838}
		% -----
			\label{EMS-1838-JS-17-SEP-1838}
			% \label{EMS-1830-JS-DAY-3LETTERMONTH-YEAR}
			
			\begin{description}
					\item[Print Sources]
				\end{description}

					\begin{tabular}{p{0.5in}p{1in}p{0.5in}p{1in}}
					JSP	 	& D6:240--245		& JSR 		& --- \\
					DC	 	& ---		& HC 		& --- \\
					HDDC 	& ---		& TCHC 		& --- \\
					\end{tabular}
					% HDDC = woodford's DC diss; JSR = Marquardt's JS Revelations; TCHC = Vogel HC
				
				\begin{description}
					\item[Online Access] \href{https://www.josephsmithpapers.org/paper-summary/letter-to-stephen-post-17-september-1838/1}{Here}
					% \item[Analysis] \hyperref[XX]{Here}
				\end{description}
				
				\begin{description}
					\item[Background]
				\end{description}
				
					% It might also be useful to give a two sentence context of the period: e.g., "This speech was given in Nauvoo to a small congregation one month prior to the destruction of the Nauvoo Expositor. These and other tensions may have been on the prophet's mind when he gave the speech."
				
				\begin{description}
					\item[Original Source]
				\end{description}
				
					% brief description of original source material, with reference to JSP or somewhere else for more info
					% Background material: it might be helpful to have a note that says whether a given document was presented as a speech, was written in a journal, etc.
					% Also a comment here about how reliable the source might be, or what shape it is in, if there are large omissions or parts that are hard to understand, and so on.
				
				\begin{description}
					\item[Text]
				\end{description}
					
					\begin{tightright}
							\hypertarget{EMS-1838-JS-17-SEP-1838-a}%
						\phantom{ }
							\marginnote{a}%
						\uline{Far West Sept 17}\supersu{th}\uline{ 1838}
					\end{tightright}
					
					Stephen Post

					Sir.

					I proceed to answer your communication of the 1\supers{st} August which I should have answered before had it not been for the press of buisness on my mind etc. The Journal is isued from this place it commenced I think in May. As to your information relative to Elders in foreign lands it is correct Elders [Orson] Hyde \& Kimble [Heber C. Kimball] with several others in company have visited Great Britton, Elders Hyde \& Kimbal have returned they are here you will probaly se[e] their narative, which I think they are about to publish in pamphlet form. They have been very successfull have baptized between one \& two thousands, ordained some 40 Elders besides other officers necessary.
						\hypertarget{EMS-1838-JS-17-SEP-1838-b}%
					Other
						\marginnote{b}%
					Stakes have been appointed but not by a committee appointed as you supposed, for the Lord has said that no stake shall be concidered a stake of Zion unless appointed, dedicated, and set apart, by the first presidency, One of these is situated about 30 miles north of this place on [Gr]and river, which is nearly as large as this place at this time \sout{One} This is called Adam Ondi Awman [Adam-ondi-Ahman], or the place where Adam dwelt. One at the mouth of Grand river called Dewitt it was a gentile city plott and named by them that is about 40 \sout{miles} miles East of this place \&c. \&c. \&c. The house of the Lord is the next question, The celler is dug the corner stones were laid July 4\supers{th} 1838. Next comes the work of the gathering. As to this, there are thousands gathering this season The road is full companies of frequently 10, 20 \& 30 <wagons> arrives, some almost daily One company which is the camp is close here with one hundred wagons John E Page report says is comming less than one hundred miles of this place, with 64 wagons and the road is litterly lined with wagons between here and Ohio.
						\hypertarget{EMS-1838-JS-17-SEP-1838-c}%
					The
						\marginnote{c}%
					work of the gathering is great. all the saints should gather as soon as possible, urge all the saints to gather immediately if they possibly can. The chance is great for purchasing lands here land is very cheap the old settlers will sell for half price yes, for quarter price they are determined to get away. Congress land is plenty and good land can be had for property other than money, such as horses wagons goods of all kinds \&c, \&c, Andiana [Indiana] and Illinois State Banks will buy Congress lands, Eastern money can be exchanged on the road, with ease. You next ask what is the cause of the papers stoping it was because the office was burnt down, by the decenters [dissenters] from the faith in Kirtland, As to Mechanical branches, all kinds are needed, \& would do well,
						\hypertarget{EMS-1838-JS-17-SEP-1838-d}%
					As
						\marginnote{d}%
					to the Stick of Joseph in the hand of Ephraim, I will merely say suppose yourself to be an Ephraimite, and suppose all this church to be, of the blood of Ephraim21 and the book of Mormon to be a record of Manasseh which would of course [be a re]cord of Joseph, Then suppose you being an Ephraimite, Should take the record of Joseph in your hand, would not then the stick of Joseph \sout{of} \sout{Joseph} be in the hand of Ephraim. solve this mistery and se[e].

					The persecutors of the saints are not asleep in Missouri but God is near as to communicate his will unto us, I can write no more at present, I would say \sout{say} may the Lord bless you and all the faithful and enable you to come up to Zion with songs of everlasting joy upon your head is the prayer of your unworthy servant and brother in the Lord Even so Amen,
					
					\begin{tightright}
					Joseph Smith Jr...
					\end{tightright}
					
		% -----
		\subsubsection{Letter to Emma Smith, 4 November 1838}
		% -----
			\label{EMS-1838-JS-4-NOV-1838}
			% \label{EMS-1830-JS-DAY-3LETTERMONTH-YEAR}
			
			\begin{description}
					\item[Print Sources]
				\end{description}

					\begin{tabular}{p{0.5in}p{1in}p{0.5in}p{1in}}
					JSP	 	& D6:279--282		& JSR 		& --- \\
					DC	 	& ---		& HC 		& --- \\
					HDDC 	& ---		& TCHC 		& --- \\
					\end{tabular}
					% HDDC = woodford's DC diss; JSR = Marquardt's JS Revelations; TCHC = Vogel HC
				
				\begin{description}
					\item[Online Access] \href{https://www.josephsmithpapers.org/paper-summary/letter-to-emma-smith-4-november-1838/1}{Here}
					% \item[Analysis] \hyperref[XX]{Here}
				\end{description}
				
				\begin{description}
					\item[Background]
				\end{description}
				
					% It might also be useful to give a two sentence context of the period: e.g., "This speech was given in Nauvoo to a small congregation one month prior to the destruction of the Nauvoo Expositor. These and other tensions may have been on the prophet's mind when he gave the speech."
				
				\begin{description}
					\item[Original Source]
				\end{description}
				
					% brief description of original source material, with reference to JSP or somewhere else for more info
					% Background material: it might be helpful to have a note that says whether a given document was presented as a speech, was written in a journal, etc.
					% Also a comment here about how reliable the source might be, or what shape it is in, if there are large omissions or parts that are hard to understand, and so on.
					
					% This one is written by JS's own hand.
				
				\begin{description}
					\item[Text]
				\end{description}
				
					\begin{tightcenter}
							\hypertarget{EMS-1838-JS-4-NOV-1838-a}%
						November
							\marginnote{a}%
						4\supers{th} 1838
					\end{tightcenter}
					
					Indipendace [Independence] Jackson Co--- Mo---
					
					My dear and beloved companion, of my bosam, in tribulation, and affliction, I woud inform you that I am well, and \sout{I am} that we are all of us in good spirits as regards our own fate, we have been protected by the Jackson County boys, in the most genteel manner, and arrived here in <the> midst of a splended perade, \sout{this} a little after noon, instead <of> going to goal [jail] we have a good house provided for us and the kind[e]st treatment, I have great anxiety about you, and my lovely children, my heart morns <and> bleeds for the brotheren, and sisters, and for the slain <of the> people of God, \sout{I} Colonal, [George M.] Hinkle, proved to be a trator, to the Church, he is worse than a hull who betraid the army at detroit, he decoyed <us> unawares God reward him, \sout{I} Johon Carl [John Corrill] \sout{told} \sout{<general Willson>} \sout{was a going} told general, [Moses] Wilson, that he was a going to leave the Church, general Willson says he thinks much less of him now then before, why I mention this is to have you careful not to trust them, if we are permited to \sout{be} stay any time here, we <have> obtained a promice that \sout{they} we may have our families brought to us,
						\hypertarget{EMS-1838-JS-4-NOV-1838-b}%
					what
						\marginnote{b}%
					God may \sout{do} do for us I do not know but I hope for the best always in all circumstances although I go unto death, I will trust in God, what outrages may be committed by the mob I know not, but expect there will be but little <or> no restraint Oh may God have mercy on us, when we arrived at the river last night an express came to gene[r]al Willson from geneal [John B.] Clark of Howard County claiming the right of command ordering us back where <or what place> God only knows, and there is some feelings betwen the offercers, I do not know where it will end, it <is> said by some that general Clark, is determined to \sout{exterminating} exterminate God has spared some of us thus far perhaps he will extend mercy in some degree toward us <yet> some of the people of this place have told me that some of the mormans may settle in this county as others <men> do \sout{the peg}
						\hypertarget{EMS-1838-JS-4-NOV-1838-c}%
					I
						\marginnote{c}%
					have some hopes that something may turn out for good to the afflicted saints, I want you to stay where you are untill you here from me again, I may send for you to bring you to me, I cannot learn much for certainty in the situation that I am in, and can only pray for deliverance, untill it is meeted out, and take every thing as it comes, with \sout{patient} patience and fortitude, I hope you will be faithful and true to every trust, I cant write much in my situation, conduct all matters as your circumstances and necesities require, may God give you wisdom and prudance and sobriety which <I> have every reason to believe you will, those little <childrens> are subjects of my meditation continually, tell them that Father is yet alive, God grant that he may see them again Oh Emma for God sake do not forsake me nor the truth but remember, if I do <not> meet you again in this life may God grant that we may <may we> meet in heaven, I cannot express my feelings, my heart is full, Farewell Oh my kind and affectionate Emma I am yours forever your Husband and true friend
					
					\begin{tightright}
					[Joseph Smith Jr.]...
					\end{tightright}
					
		% -----
		\subsubsection{Letter to Emma Smith, 12 November 1838}
		% -----
			\label{EMS-1838-JS-12-NOV-1838}
			% \label{EMS-1830-JS-DAY-3LETTERMONTH-YEAR}
			
			\begin{description}
					\item[Print Sources]
				\end{description}

					\begin{tabular}{p{0.5in}p{1in}p{0.5in}p{1in}}
					JSP	 	& D6:290--293		& JSR 		& --- \\
					DC	 	& ---		& HC 		& --- \\
					HDDC 	& ---		& TCHC 		& --- \\
					\end{tabular}
					% HDDC = woodford's DC diss; JSR = Marquardt's JS Revelations; TCHC = Vogel HC
				
				\begin{description}
					\item[Online Access] \href{https://www.josephsmithpapers.org/paper-summary/letter-to-emma-smith-12-november-1838/1}{Here}
					% \item[Analysis] \hyperref[XX]{Here}
				\end{description}
				
				\begin{description}
					\item[Background]
				\end{description}
				
					% It might also be useful to give a two sentence context of the period: e.g., "This speech was given in Nauvoo to a small congregation one month prior to the destruction of the Nauvoo Expositor. These and other tensions may have been on the prophet's mind when he gave the speech."
				
				\begin{description}
					\item[Original Source]
				\end{description}
				
					% brief description of original source material, with reference to JSP or somewhere else for more info
					% Background material: it might be helpful to have a note that says whether a given document was presented as a speech, was written in a journal, etc.
					% Also a comment here about how reliable the source might be, or what shape it is in, if there are large omissions or parts that are hard to understand, and so on.
					
					% This one is written by JS's own hand.
				
				\begin{description}
					\item[Text]
				\end{description}
				
						\hypertarget{EMS-1838-JS-12-NOV-1838-a}%
					November
						\marginnote{a}%
					12\supers{th} 1838 Richmond

					My Dear Emma.

					we are prisoners in chains, and under strong guards, for Christ sake and for no other causes although there has been things that were unbeknown to us, and altogether beyond our controal, that might seem, to the mob to be a pretext, for them to persacute us, but on examination, I think that the authorities, will discover our inocence, and set us free, but if this blessing cannot be \sout{done} obtained, I have this consolation that I am an innocent man, let what will befall me, I recieved your letter which I read over and over again, it was a sweet morsal to me,
						\hypertarget{EMS-1838-JS-12-NOV-1838-b}%
					Oh
						\marginnote{b}%
					God grant that I may have the privaliege of seeing once more my lovely Family, in the injoyment, of the sweets of liberty, and sotiaial life, to press them to my bosam and kiss\sout{ng} their lovely cheeks would fill my heart with unspeakable \sout{great} grattitude, tell the chilldren that I am alive and trust I shall come and see them before long, comfort their hearts all you can, and try to be comforted yourself, all you can, there is no possible dainger but what we shall be set at Liberty if Justice can be \sout{<and>} done <and> that you know as well as myself, the tryal will begin today for some of us, Lawyer Rice [Amos Rees] and we expect [Alexander] Doniphan, will plead our cause, we could <git> no others in time for the tryal, they are able man and <will> do well no doubt, Brother Robison [George W. Robinson] is chained next to me he \sout{he} has a true heart and a firm mind, Brother Whight [Lyman Wight], is next, Br. [Sidney] Rigdon, next, Hyram [Hyrum Smith], next, Parely [Parley P. Pratt], next, Amasa [Lyman], next, and thus we are bound together in chains as well as the cords of everlasting love,
						\hypertarget{EMS-1838-JS-12-NOV-1838-c}%
					we
						\marginnote{c}%
					are in good spirits and rejoice that we are counted worthy to be persicuted for christ sake, tell little Joseph, he must be a good boy, and Father loves him <with> a perfect love, he is the Eldest must not hurt those that <are> smaller then him, but cumfort them tell little Frederick, Father, loves him, with all his heart, he is a lovely boy Julia is a lovely little girl, I love hir also She is a promising child, tell her Father wants her to remember him and be a good girl, tell all the rest that I think of them and pray for them all, Br Babbit is waitting to carry our letters, for us \sout{the} colonal <price [Sterling Price]> is inspecting them therefore my time is short \sout{<the>} little \sout{baby} Elexander [Alexander] is on my mind continualy
						\hypertarget{EMS-1838-JS-12-NOV-1838-d}%
					Oh
						\marginnote{d}%
					my affectionate Emma, I want you to remember that I am <a> true and faithful friend, to you and the chilldren, forever, my heart is intwined around you[r]s forever and ever, Oh may God bless you all amen \sout{you} I am your husband and am in bonds and tribulation \&c
					
					\begin{tabular}{p{12cm}p{3cm}}
					to Emma Smith\} & Joseph Smith Jr
					\end{tabular}

					P S write as often as you can, and if possible come and see me, and bring the chilldren if possible, act according to your own feelings <and> best Judgement, and indeavour to be comforted, if possible, and I trust that all will turn out for the bist. yours, J.S...
					
		% -----
		\subsubsection{Letter to the Church in Caldwell County, Missouri, 16 December 1838}
		% -----
			\label{EMS-1838-JS-16-DEC-1838}
			% \label{EMS-1830-JS-DAY-3LETTERMONTH-YEAR}
			
			\begin{description}
					\item[Print Sources]
				\end{description}

					\begin{tabular}{p{0.5in}p{1in}p{0.5in}p{1in}}
					JSP	 	& D6:294--310		& JSR 		& --- \\
					DC	 	& ---		& HC 		& 3:226--233 \\
					HDDC 	& ---		& TCHC 		& 3:209--217 \\
					TPJS 	& 122--129			& 	 		&  \\
					\end{tabular}
					% HDDC = woodford's DC diss; JSR = Marquardt's JS Revelations; TCHC = Vogel HC
				
				\begin{description}
					\item[Online Access] \href{https://www.josephsmithpapers.org/paper-summary/letter-to-the-church-in-caldwell-county-missouri-16-december-1838/1}{Here}
					% \item[Analysis] \hyperref[XX]{Here}
				\end{description}
				
				\begin{description}
					\item[Background]
				\end{description}
				
					% It might also be useful to give a two sentence context of the period: e.g., "This speech was given in Nauvoo to a small congregation one month prior to the destruction of the Nauvoo Expositor. These and other tensions may have been on the prophet's mind when he gave the speech."
				
				\begin{description}
					\item[Original Source]
				\end{description}
				
					% brief description of original source material, with reference to JSP or somewhere else for more info
					% Background material: it might be helpful to have a note that says whether a given document was presented as a speech, was written in a journal, etc.
					% Also a comment here about how reliable the source might be, or what shape it is in, if there are large omissions or parts that are hard to understand, and so on.
				
				\begin{description}
					\item[Text]
				\end{description}
				
					\begin{tightright}
							\hypertarget{EMS-1838-JS-16-DEC-1838-a}%
						Liberty
							\marginnote{a}%
						Jail Missouri. Dec 16\supers{th} 1838
					\end{tightright}
					
					To the church of latter day saints in Caldwell county and the saints scattered abroad and are persecuted and made desolate and are afflicted in divers manners for christ's sake and the gospel's, and whose perils are greatly augmented by the wickedness and corruption of false brethren. May grace, mercy, and peace, be and abide with you and notwithstanding all your sufferings we assure you that you have our prayers and fervent desires for your welfare both day and night. We believe that, that God who sees us in this solitary place will hear our prayers \& reward you openly. Know assuredly dear brethren that it is for the testimony of Jesus that we are in bonds and in prison.
						\hypertarget{EMS-1838-JS-16-DEC-1838-b}%
					But
						\marginnote{b}%
					we say unto you that we consider our condition better, (notwithstanding our suffering) than those who have persecuted us and smitten us and <borne> \sout{bear} false witness against us, and we most assuredly believe that those who bear false witness against us <do> seem to have a great triumph over us for the present. But we want you to remember Haman and Mordecai you know that Haman could not be satisfied so long as he saw Mordecai at the king's gate, and he sought the life of Mordecai and the people of the jews. But the Lord so ordered that Haman was hanged upon his own gallows. So shall it come to pass with poor Haman in the last days.
						\hypertarget{EMS-1838-JS-16-DEC-1838-c}%
					Those
						\marginnote{c}%
					who have sought by their unbelief and wickedness and by the principle of mobocracy to destroy us and the people of God by killing and scattering them abroad and wilfully and maliciously delivering us into the hands of murderers desiring us to be put to death thereby having us dragged about in chains and cast into prison, and for what cause; it is because we were honest men and were determined to defend the lives of the saints at the expense of our own. I say unto you that those who have thus vilely treated us like Haman shall be hanged upon their own gallows, or in other words shall fall into their own gin and trap and ditch which they have prepared for us and shall go backward and stumble and fall, and their names shall be blotted out, and God shall reward them according to all their abominations.
						\hypertarget{EMS-1838-JS-16-DEC-1838-d}%
					Dear
						\marginnote{d}%
					brethren do not think that our hearts faint as though some strange thing had happened unto us for we have seen and been assured of all these things beforehand, and have an assurance of a better hope than that of our persecutors. Therefore God has made our shoulders broad that we can bear it. We glory in our tribulation because we know that God is with us, that he is our friend and that he will save our souls. We do not care for those that kill the body they cannot harm our souls; we ask no favors at the hands of mobs nor of the world, nor of the devil nor of his emissaries the dissenters. We have never dissembled nor will we for the sake of our lives. Forasmuch then as we know that we have been endeavoring with all our mights, minds, and strength to do the will of God and all things whatsoever he has commanded us. And as to our light speeches from time to time they have nothing to do with the fixed principle of our hearts. Therefore it sufficeth us to say that our souls were vexed from day to day.
						\hypertarget{EMS-1838-JS-16-DEC-1838-e}%
					We
						\marginnote{e}%
					refer you to Isaiah who considers those who make a man an offender for a word and lay a snare for them that reproveth in the gate. We believe the old prophet verily told the truth. We have no retraction to make, we have reproved in the gate and men have laid snares for us we have spoken words and men have made us offenders, and notwithstanding all this our minds are not darkened but feel strong in the Lord. But behold the words of the savior, if the light which is in you become darkness behold how great is that darkness. Look at the dissenters. And again if you were of the world the world would love its own Look at M\supers{r} [George M.] Hinkle. A wolf in sheep's clothing. Look at his brother John Corrill Look at the beloved brother Reed Peck who aided him in leading us, as the savior was led, into the camp as a lamb prepared for the slaughter and a sheep dumb before his shearer so we opened not our mouth But these men like Balaam being greedy for a reward sold us into the hands of those who loved them, for the world loves his own.
						\hypertarget{EMS-1838-JS-16-DEC-1838-f}%
					I
						\marginnote{f}%
					would remember W[illiam] W. Phelps who comes up before us as one of Job's comforters. God suffered such kind of beings to afflict Job, but it never entered into their hearts that Job would get out of it all. This poor man who professes to be much of a prophet has no other dumb ass to ride but David Whitmer to forbid his madness when he goes up to curse Israel, and this ass not being of the same kind of Balaams therefore the angel notwithstanding appeared unto him yet he could not penetrate his understanding sufficiently so but what he brays out cursings instead of blessings. Poor ass whoever lives to see it will see him and his rider perish like those who perished in the gainsaying of Core, or after the same condemnation. Now as for these and the rest of their company we will not presume to say that the world loves them but we presume to say that they love the world and we classify them in the error Balaam and in the gainsaying of Core and with the company of Cora [Korah] and Dathan and Abiram.
						\hypertarget{EMS-1838-JS-16-DEC-1838-g}%
					Perhaps
						\marginnote{g}%
					our brethren may say because we thus write that we are offended at those characters, if we are, it is not for a word neither because they reproved in the gate. But because they have been the means of shedding innocent blood. Are they not murderers then at heart? Are not their consciences seared as with a hot iron? We confess that we are offended but the saviour said that offences must come but woe unto them by whom they come, and again blessed are ye when all men shall revile you and speak all manner of evil against you falsely for my sake, rejoice and be exceeding glad for great is your reward in heaven for so persecuted they the prophets which were before you. Now dear brethren if any men ever had reason to claim this promise we are the men, for we know that the world not only hates us but \sout{but} speak all manner of evil of us falsely for no other reason than because we have been endeavoring to teach the fulness of the gospel of Jesus Christ after we were bartered away by Hinkle and were taken into the militia camp we had all the evidence we could have wished for that the world hated us and that most cordially too.
						\hypertarget{EMS-1838-JS-16-DEC-1838-h}%
					If
						\marginnote{h}%
					there were priests of all the different sects they hated us, if there were Generals they hated us, if there were Colonels they hated us, and the soldiers and officers of every kind hated us, and the most profane blasphemers and drunkards \& whoremongers hated us, they all hated us most cordially. And now what did they hate us for, purely because of the testimony of Jesus Christ. Was it because we were liars? We know that it has been reported by some but it has been reported falsely Was it because we have committed treason against the government in Daviess County or of burglary, or of larceny or arson, or any other unlawful act in Daviess county. We know that certain priests and certain lawyers and certain judges who are the instigators aiders and abettors of a certain gang of murderers and robbers who have been carrying on a scheme of mobocracy to uphold their priestcraft against the saints of the last days for a number of years and have tried by a well contemplated and premeditated scheme to put down by physical power a system of relig[i]on that all the world by all their mutual attainments and by any fair means whatever were not able to resist.
						\hypertarget{EMS-1838-JS-16-DEC-1838-i}%
					Hence,
						\marginnote{i}%
					mobbers were encouraged by priests and Levites, by the Pharisees, Sadducees, and Essenees, and the Herodians, and the most ruthless, abandoned, and debauched, lawless inhuman and the most beastly set of men that the earth can boast of; and indeed a parallel cannot be found any where else; to gather together to steal to plunder to starve and to exterminate and burn the houses of the Mormons these are the characters that by their treasonable and avert acts have desolated and laid waste Daviess County these are the characters that would fain make all the world believe that we are guilty of the above named acts. But they represent us falsely; we say unto you that we have not committed treason, nor any other unlawful act in Daviess County was it for murder in Ray county against mob-militia who was a wolf in the first instance hide and Hair, teeth, and legs, and tail, who afterwards put on a militia sheepskin with the wool on, who can sally forth in the day time into the flock and snarl \& show his teeth, and scatter and devour the flock and satiate himself upon his prey, and then sneak back into the brambles in order that he might conceal himself in his well tried skin with the wool on.
						\hypertarget{EMS-1838-JS-16-DEC-1838-j}%
					We
						\marginnote{j}%
					are well aware that there is a certain set of priests \& satellites and mobbers that would fain make all the world believe that we are the dogs that barked at this howling wolf that made such havoc among the sheep who when he retreated howled and bleated at such a desperate rate that if one could have been there he would have thought that all the wolves whether wrapped up in sheep skins or goat skins or any other skins and in fine all the beast of the forest were awfully alarmed and catching the scent of innocent blood they sallied forth with a tremenduous howl and crying of all sorts and such a howling and such a tremenduous havoc never was known such a piece of inhumanity and relentless cruelty and barbarity cannot be found in all the annals of history. These are the characters that would make the world believe that we had committed murder by making an attack upon this howling wolf while we were at home and in our beds and asleep and knew nothing of that transaction any more than we know what is going on in China while we are within these walls.
						\hypertarget{EMS-1838-JS-16-DEC-1838-k}%
					Therefore
						\marginnote{k}%
					we say again unto you we are innocent of these things they have represented us falsely Was it for committing adultery, we are aware that false slander has gone abroad for it has been reiterated in our <ears>. These are falsehoods also. Renegadoes, mormon dissenters are running through the world and spreading various foul and libelous reports against us thinking thereby to gain the friendship of the world because they know that we are not of the world and that the world hates us; therefore they make a tool of these fellows by them they do all the injury they can and after that they hate them worse than they do us because they find them to be base traitors and sycophants. \sout{God} Such characters God hates we cannot love them the world hates them and we sometimes think the devil ought to be ashamed of them. We have heard that it has been reported by some that some of us should have said that we not only dedicated our property but our families also to the Lord, and satan taking advantage of this has transfigured it into lasciviousness such as a community of wives which is an abomination in the sight of God.
						\hypertarget{EMS-1838-JS-16-DEC-1838-l}%
					When
						\marginnote{l}%
					we consecrate our property to the Lord it is to administer to the wants of the poor and needy for this is the law of God it is not for the purpose of the rich those who have no need and when a man consecrates or dedicates his wife and children he does not give them to his brother or to his neighbor for there is no such law for the law of God is thou shalt not commit adultery thou shalt not covet thy neighbor's wife. He that looketh upon a woman to lust after her has committed adultery already in his heart. Now for a man to consecrate his property and his wife \& children to the Lord, is nothing more nor less than to feed the hungry, clothe the naked, visit the widow and the fatherless, the sick, and the afflicted, and do all he can to administer to their relief in their afflictions, and for him and his house to serve the Lord. In order to do this he and all his house must be virtuous and shun every appearance of evil. Now if any person has represented any thing other wise than what we now write he or she is a liar and have represented us falsely. And this is another manner \sout{of} of evil which is spoken against us falsely.
						\hypertarget{EMS-1838-JS-16-DEC-1838-m}%
					We
						\marginnote{m}%
					have learned also since we have been in prison that many false and pernicious things which were calculated to lead the saints far astray and to do great injury <have been taught by Dr. [Sampson] Avard> as coming from the Presidency \sout{taught by Dr Avard} and we have reason to fear <that> many <other \sout{things}> designing and corrupt characters like unto himself <have been teaching many things> which the presidency never knew of being taught in the church by any body untill after they were made prisoners, which if they had known of, they would have spurned them and their authors from them as they would the gates of hell. Thus we find that there has been frauds and secret abominations and evil works of darkness going on leading the minds of the weak and unwary into confusion and distraction, and palming it all the time upon the presidency while mean time the presidency were ignorant as well as innocent of these things, which were practicing in the church in their name and were attending to their own family concerns, weighed down with sorrow, in debt, in poverty, in hunger assaying to be fed yet finding themselves receiving deeds of charity but inadequate to their subsistence, and because they received those deeds they were envied and hated by those who professed to be their friends
						\hypertarget{EMS-1838-JS-16-DEC-1838-n}%
					But
						\marginnote{n}%
					notwithstanding we thus speak we honor the church when we speak of the church, as a church, for their liberality, kindness, patience, and long suffering, and their continued kindness towards us. And now brethren we say unto you, what can we enumerate more; is not all manner of evil of every description spoken against us falsely, yea, we say unto unto you falsely; we have been misrepresented and misunderstood and belied and the purity of our hearts have not been known. And it is through ignorance, yea, the very depth of ignorance is the cause of it, and not only ignorance but gross wickedness on the part of some and hypocrisy also who by a long face and sanctified prayers and very pious sermons had power to lead the minds of the ignorant and unwary and thereby obtain such influence that when we approached their iniquities the devil gained great advantage \& would bring great sorrow upon our heads and in fine we have waded through an ocean of tribulation, and mean abuse practiced upon us by the ill bred and ignorant such as Hinkle, Corrill, and Phelps, Avard, Reed Peck, [John] Cleminson, and various others who are so very ignorant that they cannot appear respectable in any decent and civilized society, and whose eyes are full of adultery and cannot cease from sin. Such characters as [William E.] McLellin, John Whitmer, D. Whitmer, O[liver] Cowdery, Martin Harris, who are too mean to mention and we had liked to have forgotten them.
						\hypertarget{EMS-1838-JS-16-DEC-1838-o}%
					[Thomas B.]
						\marginnote{o}%
					Marsh \& [Orson] Hyde whose hearts are full of corruption, whose cloak of hypocrisy was not sufficient to shield them or to hold them up in the hour of trouble, who after having escaped the pollutions of the world through the knowledge of God and become again entangled and overcome the latter end is worse than the first. But it has happened unto them according to the words of the savior, the dog has returned to his vomit, and the sow that was washed to her wallowing in the mire. Again if we sin wilfully after we have received the knowledge of the truth, there remaineth no more sacrifice for sin, but a certain fearful looking <for> of judgement and fiery indignation to come which shall devour these adversaries. For he who despiseth Moses' law died without mercy under two or three witnesses of how much more severe punishment suppose ye shall he be thought worthy who hath sold his brother and denied the new and everlasting covenant by which he was sanctified calling it an unholy thing and doing despite to the spirit of grace.
						\hypertarget{EMS-1838-JS-16-DEC-1838-p}%
					And
						\marginnote{p}%
					again we say unto you that inasmuch as there be virtue in us and the holy priesthood hath been conferred upon us, and the keys of the kingdom hath not been taken from us, for verily thus saith the Lord be of good cheer for the keys that I gave unto <you> are yet with you Therefore we say unto you dear brethren in the name of the Lord Jesus Christ, we deliver these characters unto the buffetings of satan untill the day of redemption that they may be dealt with according to their works and from henceforth their works shall be made manifest. And now dear and well beloved brethren and when we say brethren we mean those who have continued faithful in christ men, women, and children, we feel to exhort you in the name of the Lord Jesus, to be strong in the faith of the new and everlasting covenant, and nothing frightened at your enemies. For what has happened unto us is an evident token to them of damnation but unto us of salvation and that of God. Therefore hold on even unto death, for he that seeks to save his life shall \sout{loose} lose it but he that loseth his life for my sake and the gospels shall find it sayeth Jesus Christ.
						\hypertarget{EMS-1838-JS-16-DEC-1838-q}%
					Brethren
						\marginnote{q}%
					from henceforth let truth and righteousness prevail and abound in you and in all things be temperate, abstain from every appearance of evil, drunkenness, and profane language, and from every thing which is unrighteous or unholy; also from enmity, and hatred, and covetousness and from every unholy desires. Be honest one with another, for it seemeth that some have come short of these things, and some have been uncharitable \& have manifested greediness because of their debts towards those who have been persecuted \& dragged about with chains without cause and imprisoned. Such persons God hates and they shall have their turn of sorrow in the rolling of the great wheel for it rolleth and none can hinder. Zion shall yet live though she seemeth to be dead.
						\hypertarget{EMS-1838-JS-16-DEC-1838-r}%
					Remember
						\marginnote{r}%
					that whatsoever measure you meet out to others it shall be measured to you again. We say unto you brethren be not afraid of your adversaries contend earnestly against mobs, and the unlawful works of dissenters and of darkness. And the very God of peace shall be with you and make a way for your escape from the adversary of your souls we commend you to God and the word of his grace which is able to make us wise unto salvation. Amen.
					
					\begin{tightright}
						Joseph Smith Jun.
					\end{tightright}
					
					[\textit{page [8] blank}]
			
% -----------------------------------------------------------
\pagebreak{}
% -----------------------------------------------------------